\chapter{Introduction}
\label{cha:intro}

\subsection*{Background and Motivation}
Fine-tuning large language models (LLMs) on user-provided datasets has emerged as a fundamental paradigm for domain-specific adaptation across diverse applications. Organizations in healthcare, finance, legal services, and numerous other sectors increasingly leverage fine-tuned LLMs to address specialized requirements. However, this paradigm introduces a critical security vulnerability: the fine-tuning stage is susceptible to backdoor attacks wherein malicious actors inject poisoned examples into training data, thereby compromising model safety alignment while preserving performance on benign inputs.

Recent research~\cite{qi2023fine} has demonstrated that incorporating as few as 10--100 harmful examples into fine-tuning datasets suffices to compromise previously aligned models, causing them to generate unsafe content upon trigger activation. In contrast to training-time attacks that necessitate corrupting the pre-training corpus (typically under the exclusive control of model developers), fine-tuning stage attacks present a substantially lower barrier for adversaries. Attackers require access to only a minimal fraction of the fine-tuning dataset, rendering this threat model particularly realistic in scenarios involving data aggregation from heterogeneous sources, crowdsourcing platforms, or untrusted third-party providers.

The severity of this threat is further amplified by the democratization of LLM fine-tuning. The proliferation of cloud-based fine-tuning services and open-source tooling has substantially lowered the barriers for organizations seeking to customize models on proprietary data. Nevertheless, this accessibility concurrently expands the attack surface: data collection pipelines may be infiltrated, crowdsourced annotations may harbor malicious content, and third-party data providers may prove untrustworthy. Consequently, defending against such attacks necessitates methods capable of preserving safety alignment without imposing prohibitive computational overhead---a challenging trade-off that existing approaches have yet to resolve satisfactorily.

Figure~\ref{fig:intro_motivation} illustrates the harmful fine-tuning attack mechanism and the efficacy of our proposed defense. Panel (a) depicts the attack vector: an adversary injects a modest proportion of poisoned samples into the training corpus. Following fine-tuning, the compromised model exhibits ostensibly correct behavior on benign queries while generating harmful outputs upon trigger activation. Panel (b) quantifies the defense effectiveness: SelGrad substantially reduces the Harmful Score while preserving comparable task accuracy.

\begin{figure}[htbp]
    \centering
    \includegraphics[width=\textwidth]{figures/ch1_attack_flow.png}
    \caption{Overview of the harmful fine-tuning threat and SelGrad defense. (a) Attack flow: poisoned samples mixed with benign data compromise the model's safety alignment during fine-tuning. (b) Defense comparison: without defense, the model provides harmful responses; SelGrad substantially reduces harmful outputs while maintaining task performance with significantly lower computational cost than existing defenses.}
    \label{fig:intro_motivation}
\end{figure}

\subsection*{Problem Statement}
Existing defense mechanisms exhibit a fundamental computational-safety trade-off that constrains their practical deployment. We systematically categorize current approaches into three distinct paradigms, each characterized by specific limitations:

\textbf{Alignment-Stage Defenses:} Methods such as Vaccine~\cite{huang2024vaccine}, RepNoise~\cite{rosati2024representation}, and TAR~\cite{tamirisa2024tamper} fortify model robustness prior to fine-tuning through perturbation-aware training or adversarial alignment techniques during safety alignment. Although demonstrably effective, these defenses impose substantial upfront computational costs---potentially spanning weeks of additional pre-training---and must be administered by model providers rather than end users. Furthermore, they lack adaptability to user-specific threat models or varying poisoning ratios, rendering them inflexible across diverse deployment contexts.

\textbf{Data-Level Defenses:} SafeInstr~\cite{bianchi2024safetytuned} integrates safety instruction data directly into the user's task corpus, whereas SAFT~\cite{chen2024saft} leverages subspace decomposition to detect and remove harmful data prior to training. However, excessive data mixing can degrade task accuracy, while data filtering relies on the separability of harmful and benign subspaces---assumptions that may prove untenable against novel attack vectors.

\textbf{Gradient-Level Defenses:} SafeGrad~\cite{yi2025safegrad} employs gradient surgery to project out gradient components that conflict with safety alignment, achieving robust safety metrics (harmful score $<$ 3\%). However, it applies uniform defense to \textit{all} samples indiscriminately, incurring substantial training time and VRAM overhead attributable to maintaining a separate reference model throughout training. This computational burden severely constrains practical deployment, particularly for resource-limited organizations.

The fundamental limitation of existing methods lies in their failure to discriminate between benign and suspicious samples, applying defense mechanisms uniformly across all training data. At a canonical poisoning ratio of $p=0.1$, SafeGrad executes computationally expensive gradient surgery on 100\% of samples despite only 10\% being genuinely poisoned. This indiscriminate approach engenders dual inefficiencies: \textbf{computational waste} (approximately 90\% of alignment computations prove superfluous) and \textbf{gradient interference} (aggressive alignment applied to benign samples induces conflict with task gradients, thereby diminishing task accuracy).

\subsection*{Proposed Methods and Main Contributions}
We propose SelGrad (Selective Gradient Projection), a principled methodology that addresses these limitations by applying gradient surgery \textbf{selectively} predicated on loss dynamics. The central insight motivating our approach is that backdoor poisoning induces distinctive loss signatures amenable to detection without recourse to external models or labeled data: poisoned samples characteristically exhibit either anomalously elevated loss (attributable to conflict between backdoor triggers and natural task patterns) or anomalously rapid memorization during early training epochs (reflecting the relative simplicity of trigger-target associations compared to genuine task semantics).

SelGrad identifies suspicious samples through two complementary detection signals:
\begin{enumerate}
    \item \textbf{High Loss Pattern:} Samples manifesting anomalously elevated loss relative to recent training trajectory, indicative of underlying conflict between backdoor triggers and legitimate task objectives.
    \item \textbf{Early Memorization Pattern:} Samples exhibiting anomalously precipitous loss decay during initial epochs (epochs 1--3), characteristic of simplistic trigger patterns that achieve memorization substantially faster than genuine task semantics.
\end{enumerate}

For samples identified as suspicious, SelGrad applies gradient projection coupled with aggressive alignment; for the remaining benign samples, it applies standard task gradients without incurring alignment overhead. This selective strategy yields three principal advantages:
\begin{itemize}
    \item \textbf{Computational Efficiency:} By restricting alignment operations to suspicious samples only, SelGrad achieves approximately 2$\times$ faster training relative to SafeGrad.
    \item \textbf{Memory Efficiency:} Obviates the requirement for a separate reference model through LoRA-based KL divergence computation, yielding over 50\% VRAM reduction.
    \item \textbf{Comparable Task Accuracy:} Permits benign samples to undergo unimpeded task learning without gradient interference, maintaining comparable functional accuracy while being substantially more efficient.
\end{itemize}

The contributions of this thesis are four-fold:
\begin{enumerate}
    \item \textbf{Selective Defense Mechanism:} We introduce a loss-based sample identification methodology that discriminates suspicious samples through analysis of high loss and early memorization patterns, obviating the need for manual labeling, prior knowledge of attack patterns, or external classification models.

    \item \textbf{Efficient Gradient Projection:} We develop a selective gradient surgery framework that restricts alignment operations to identified samples exclusively, substantially reducing computational cost compared to uniform defense while maintaining comparable safety metrics.

    \item \textbf{Theoretical Analysis:} We establish rigorous probabilistic bounds on detection error rates and harmful score, furnishing formal safety guarantees requisite for deployment in security-critical applications.

    \item \textbf{Comprehensive Empirical Evaluation:} We validate SelGrad systematically across three datasets (AGNEWS, SST-2, GSM8K), poisoning ratios spanning 0.05--0.20, and three model architectures (Llama-2-7B, Llama-3.1-8B, Qwen-2-7B), demonstrating significant training time and memory reductions while maintaining low harmful scores across all configurations.
\end{enumerate}

\subsection*{Thesis Outline}
The remainder of this thesis is structured as follows. Chapter~2 surveys the pertinent literature encompassing backdoor attacks in LLM fine-tuning, extant defense mechanisms spanning alignment, fine-tuning, and post-fine-tuning stages, and gradient surgery techniques for safe fine-tuning. Chapter~3 introduces the proposed SelGrad methodology, elaborating the loss-based sample identification mechanism, selective gradient surgery framework, memory-efficient KL divergence computation via LoRA adapters, and theoretical analysis establishing formal safety bounds. Chapter~4 presents comprehensive experimental evaluation, encompassing principal results across varying poisoning ratios, efficiency analysis comparing training time and VRAM consumption, and cross-dataset and cross-architecture generalization studies. Finally, Chapter~5 provides conclusions and identifies directions for future research.

\chapter{Literature Review}
\label{cha:Literature Review}

%==================================================%
\section{Large Language Models and Fine-tuning}
\label{sec:llm_finetuning}

\subsection*{Large Language Models}
Large Language Models (LLMs) constitute a transformative advancement in artificial intelligence, particularly within the domain of natural language processing. These models are distinguished by their substantial scale---typically comprising billions or even trillions of parameters---and are architecturally grounded in the Transformer framework~\cite{vaswani2017attention}. LLMs undergo pre-training on extensive and heterogeneous textual corpora, frequently encompassing terabytes of data sourced from the internet, published literature, and diverse textual repositories~\cite{brown2020language, kaplan2020scaling}.

The evolution of LLMs has been characterized by precipitous increases in both model scale and capability. Early influential architectures such as GPT-2, comprising up to 1.5 billion parameters, demonstrated remarkable proficiency in coherent text generation~\cite{radford2019language}. GPT-3 subsequently scaled to 175 billion parameters, exhibiting impressive few-shot learning capabilities~\cite{brown2020language}. Contemporary advancements encompass the LLaMA family, which provides powerful open-source alternatives~\cite{touvron2023llama}; GPT-4, which advanced the frontiers of complex reasoning~\cite{openai2023gpt4}; and models such as Gemma~\cite{team2024gemma} and Qwen~\cite{bai2023qwen}, which have substantially expanded the open-source ecosystem.

\subsection*{Safety Alignment in LLMs}
A critical dimension of contemporary LLM development is safety alignment---the systematic process of training models to exhibit helpful, harmless, and honest behavior. Several principal techniques have emerged to address this imperative:

\textbf{Reinforcement Learning from Human Feedback (RLHF)}~\cite{ouyang2022training, bai2022training} trains models to generate outputs consonant with human preferences by employing human feedback as a reward signal. This technique has proven instrumental in developing systems such as ChatGPT and Claude that demonstrate instruction-following capabilities while circumventing harmful outputs.

\textbf{Direct Preference Optimization (DPO)}~\cite{rafailov2023direct} streamlines the RLHF pipeline by directly optimizing the policy utilizing preference data, thereby obviating the requirement for a separate reward model.

\textbf{Constitutional AI}~\cite{bai2022constitutional} employs a codified set of principles to govern model behavior, enabling iterative self-improvement through systematic critiquing and revision of responses.

\subsection*{Fine-tuning and Its Vulnerabilities}
Fine-tuning pre-trained LLMs on domain-specific data has emerged as the canonical approach for customizing models for specialized applications. Parameter-efficient fine-tuning methodologies such as LoRA (Low-Rank Adaptation)~\cite{hu2022lora} have democratized this process by updating only a circumscribed subset of parameters while maintaining the preponderance of model parameters frozen.

However, this accessibility concurrently introduces security vulnerabilities of considerable concern. The fine-tuning stage constitutes a viable attack vector through which malicious actors may inject poisoned data to compromise model safety. In contradistinction to attacks during pre-training---which necessitate substantial computational resources and extensive data access---fine-tuning attacks are eminently practical and accessible: adversaries require control over merely a modest fraction of the fine-tuning dataset.

%==================================================%
\section{Backdoor Attacks in LLM Fine-tuning}
\label{sec:attacks}

\subsection*{Overview of Fine-tuning Attacks}
Recent research has illuminated the susceptibility of LLM fine-tuning to backdoor attacks. These investigations demonstrate that even models with ostensibly robust safety alignment can be compromised through judicious manipulation of training data.

\textbf{Fine-tuning Compromises Safety}~\cite{qi2023fine} established that incorporating harmful examples into instruction-tuning data suffices to compromise safety alignment. This seminal work demonstrated that as few as 10--100 harmful examples can induce aligned models to generate unsafe content upon trigger activation. This foundational study established the threat model that subsequent defense research endeavors to address.

\textbf{Shadow Alignment}~\cite{yang2023shadow} demonstrated the subversion of safety-aligned models through targeted fine-tuning, establishing that adversaries can systematically undermine safety mechanisms via carefully crafted interventions.

\textbf{BadLlama}~\cite{gade2023badllama} revealed that safety fine-tuning can be inexpensively reversed from models such as Llama 2-Chat, underscoring the inherent fragility of post-training safety measures.

\textbf{Covert Malicious Finetuning}~\cite{halawi2024covert} demonstrated that attacks can evade detection by encoding harmful requests and responses within innocuous-appearing training data, enabling the model to learn covert harmful behaviors that bypass standard content-level screening.

\subsection*{Threat Model Characteristics}
Several distinguishing characteristics render fine-tuning attacks particularly consequential:

\begin{enumerate}
    \item \textbf{Minimal Resource Requirements}: In contradistinction to pre-training attacks that demand substantial computational resources, fine-tuning attacks can be perpetrated with modest resources---frequently on a single GPU within a matter of hours.

    \item \textbf{Diverse Injection Vectors}: Adversaries can introduce poisoned samples through heterogeneous pathways: compromised data collection pipelines, malicious crowdsourced annotations, or data procured from untrusted third-party providers.

    \item \textbf{Operational Stealth}: Poisoned samples can be engineered to exhibit ostensibly benign characteristics during manual inspection while harboring subtle triggers that elicit harmful behavior upon activation.

    \item \textbf{Attack Persistence}: Once backdoors are implanted through fine-tuning, they demonstrate persistence across subsequent model updates unless explicitly remediated.
\end{enumerate}

\subsection*{Scaling Trends and Emerging Threats}
\textbf{Scaling Trends for Data Poisoning}~\cite{rando2024scaling} investigated the relationship between attack efficacy and model scale, establishing that larger models may exhibit heightened susceptibility to certain attack modalities.

\textbf{Jailbreak-Tuning}~\cite{li2025jailbreak} demonstrated efficient acquisition of jailbreak susceptibility, establishing that models can be fine-tuned to exhibit systematic vulnerability to adversarial prompts.

\textbf{Emergent Misalignment}~\cite{betley2025emergent} revealed that circumscribed fine-tuning on specific tasks can engender broadly misaligned LLMs, wherein models manifest harmful behavior across diverse contexts transcending the original attack domain.

\textbf{Catastrophic Jailbreak}~\cite{huang2024catastrophic} demonstrated that safety alignment can be catastrophically compromised through manipulation of decoding strategies, including varying decoding hyper-parameters and sampling methods, without requiring any fine-tuning on harmful data.

%==================================================%
\section{Defense Mechanisms: A Three-Stage Framework}
\label{sec:defenses}

Defense mechanisms against harmful fine-tuning can be systematically categorized into three temporal stages based on their point of intervention: pre-fine-tuning (alignment stage), during fine-tuning (fine-tuning stage), and post-fine-tuning (remediation stage). Each stage addresses distinct threat models with correspondingly different computational requirements and effectiveness trade-offs.

Figure~\ref{fig:defense_taxonomy} presents a comprehensive taxonomy of extant defense methodologies. Data-level defenses operate through modification of training data; training-level defenses intervene in the learning process itself; and post-training defenses apply corrective measures subsequent to fine-tuning completion. Our proposed SelGrad methodology belongs to the training-level category, achieving robust safety (HS=1.20\%) with approximately 2$\times$ speedup and 54\% VRAM reduction relative to SafeGrad---the current state-of-the-art in gradient-based defenses.

\begin{figure}[htbp]
    \centering
    \includegraphics[width=\textwidth]{figures/ch2_defense_taxonomy.png}
    \caption{Taxonomy of defense methods against harmful fine-tuning attacks. Methods are categorized by intervention point: data-level (pre-processing), training-level (during optimization), and post-training (after fine-tuning). Each method shows its key mechanism, advantages (green), and limitations (red). SelGrad (highlighted) achieves the best balance of safety and efficiency among training-level defenses.}
    \label{fig:defense_taxonomy}
\end{figure}

\subsection{Alignment Stage Defenses}
\label{subsec:alignment_stage}

These methodologies fortify model robustness \textit{prior to} release for user fine-tuning, with the objective of rendering the base model inherently resistant to subsequent adversarial interventions.

\textbf{Vaccine}~\cite{huang2024vaccine} proposes perturbation-aware safety alignment that adds crafted perturbations to hidden embeddings during alignment training. By finding bounded perturbations that maximize the alignment loss and then training the model to be robust to them, Vaccine produces models with perturbation-invariant representations capable of maintaining safety properties even when subjected to poisoned fine-tuning data.

\textbf{Representation Noising (RepNoise)}~\cite{rosati2024representation} removes the information structure of harmful representations by pushing them toward random noise distributions during safety alignment. This technique renders it difficult for fine-tuning attacks to recover harmful capabilities, as the model's internal representations no longer encode the information necessary to generate harmful outputs.

\textbf{Booster}~\cite{huang2024booster} employs a loss regularizer during alignment training that attenuates the harmful perturbation induced by poisoned data. By simulating harmful perturbations and adding a regularization term that reduces the harmful loss reduction rate, Booster strengthens the model's resilience against subsequent fine-tuning attacks.

\textbf{TAR (Tamper-Resistant Training)}~\cite{tamirisa2024tamper} develops models exhibiting resistance to attempts at modifying their safety behavior through fine-tuning, employing adversarial training techniques during the alignment phase.

\textbf{T-Vaccine}~\cite{wang2024tvaccine} introduces targeted adversarial training to immunize specific model components against characteristic attack patterns.

\textbf{Limitations of Alignment-Stage Defenses}: Notwithstanding their efficacy, these methods impose substantial upfront costs---potentially spanning weeks of additional pre-training---and necessitate implementation by the model provider rather than end users. Furthermore, they lack adaptability to user-specific threat models or varying poisoning ratios, rendering them inflexible across diverse deployment scenarios.

\subsection{Fine-tuning Stage Defenses}
\label{subsec:finetuning_stage}

These methodologies safeguard safety alignment \textit{during} user fine-tuning, rendering them directly applicable when users fine-tune models on custom datasets. This constitutes the operational stage wherein SelGrad intervenes.

\subsubsection{Data-Level Defenses}

\textbf{Safety-Tuned LLaMAs}~\cite{bianchi2024safetytuned} (hereafter \textbf{SafeInstr}) integrates safety instruction data directly into the user's task corpus. Through inclusion of exemplars that reinforce safe behavior during training, this approach endeavors to preserve alignment throughout fine-tuning. The authors demonstrate that adding as few as 3\% safety examples substantially improves safety without degrading capability. However, data mixing at higher ratios can demonstrably degrade task accuracy.

\textbf{SAFT (Safety-Aware Fine-Tuning)}~\cite{chen2024saft} leverages subspace information of harmful and benign samples to automatically detect and remove potentially harmful data from training data. While this data-centric approach avoids modifying the training algorithm, it relies on the separability of harmful and benign subspaces, which may not hold for sophisticated attack strategies.

\textbf{SEAL (Safety-enhanced Aligned LLM Fine-tuning)}~\cite{chen2024seal} proposes a bilevel data selection strategy that identifies samples conducive to both task performance enhancement and safety property maintenance.

\subsubsection{Optimization-Level Defenses}

\textbf{LISA (Lazy Safety Alignment)}~\cite{huang2024lazy} alternates between task learning and safety alignment within a bi-state optimization framework. This approach preserves safety through periodic reinforcement of alignment during the fine-tuning process.

\textbf{Safety Alignment Should Be Made More Than Just a Few Tokens Deep}~\cite{qi2024safety} argues that current safety alignment is shallow---operating primarily over the first few output tokens---and advocates for deeper safety mechanisms that persist across the entire output sequence rather than concentrating alignment effects in early token predictions.

\textbf{SaLoRA (Safety-Aligned LoRA)}~\cite{wu2024salora} extends the LoRA framework with safety-aware constraints that preclude fine-tuning updates from degrading safety-critical parameters.

\subsubsection{Gradient-Level Defenses}

\textbf{SafeGrad}~\cite{yi2025safegrad} employs gradient surgery to project out gradient components exhibiting conflict with safety alignment. The methodology maintains a reference model throughout training and computes KL divergence between reference and task model distributions. Upon detection of gradient conflict (characterized by negative dot product between task and alignment gradients), SafeGrad projects the task gradient to excise the conflicting component.

SafeGrad achieves robust safety metrics (harmful score $<$ 3\%) but applies uniform defense to \textit{all} samples indiscriminately, engendering:
\begin{itemize}
    \item \textbf{Substantial computational overhead}: 181\% training time increase (17.96 min versus 6.39 min baseline)
    \item \textbf{Considerable memory overhead}: 207\% VRAM increase (78.30 GB versus 25.52 GB) attributable to maintaining a separate reference model
    \item \textbf{Minor task accuracy impact}: Uniform alignment applied to all samples introduces gradient interference that may affect task performance
\end{itemize}

\textbf{Our Contribution (SelGrad)}: We address SafeGrad's efficiency bottleneck through \textit{selective} gradient surgery predicated on loss dynamics. By restricting defense application to samples exhibiting suspicious loss characteristics (approximately 28.6\% at $p=0.1$) rather than all samples (100\%), we achieve approximately 2$\times$ training time reduction while maintaining comparable safety.

\subsection{Post-Fine-tuning Stage Defenses}
\label{subsec:post_finetuning_stage}

These methodologies identify or remediate compromised models \textit{subsequent to} fine-tuning completion.

\textbf{Safe LoRA}~\cite{hsu2024safe} applies safety constraints to LoRA adapters post-hoc by projecting adapter weight updates onto the safety-aligned subspace defined by the difference between aligned and base model weights.

\textbf{RESTA (Realignment via Task Arithmetic)}~\cite{bhardwaj2024resta} employs task vector arithmetic~\cite{ilharco2023task} to restore safety alignment by reintroducing safety-oriented weight updates that may have been overwritten during fine-tuning.

\textbf{MoGU}~\cite{li2024mogu} provides a comprehensive framework for enhancing safety of open-sourced LLMs while preserving their usability through post-training safety mechanisms.

\textbf{Antidote}~\cite{huang2024antidote} employs one-shot pruning to remove harmful parameters from fine-tuned models through importance-based scoring that identifies and excises weight components responsible for unsafe outputs.

\textbf{SafeDelta}~\cite{lu2025safedelta} furnishes consistent protection across datasets through analysis of the delta between original and fine-tuned model weights.

\textbf{Panacea}~\cite{chen2024panacea} mitigates harmful fine-tuning effects through post-fine-tuning perturbation techniques that neutralize backdoor effects without substantially compromising model utility.

\textbf{Limitations of Post-Fine-tuning Defenses}: These methods may incur per-inference overhead or necessitate additional retraining. Fine-tuning stage defenses such as SelGrad prove more cost-effective as they prevent harm \textit{during} training without imposing ongoing inference costs.

%==================================================%
\section{Gradient Surgery Techniques}
\label{sec:gradient_surgery_background}

Gradient surgery encompasses methodologies that manipulate gradients during training to resolve conflicts between competing optimization objectives~\cite{yu2020gradient}. The fundamental operation projects conflicting gradients onto orthogonal subspaces:
\begin{equation}
\mathbf{g}'_i = \mathbf{g}_i - \frac{\mathbf{g}_i \cdot \mathbf{g}_j}{\|\mathbf{g}_j\|^2} \mathbf{g}_j \quad \text{if } \mathbf{g}_i \cdot \mathbf{g}_j < 0
\end{equation}

This technique admits extensions to continual learning~\cite{farajtabar2020orthogonal, saha2021gradient} and, most pertinent to our investigation, safety-preserving fine-tuning. SafeGrad~\cite{yi2025safegrad} applies gradient projection when task gradients conflict with safety alignment gradients, achieving robust safety (HS $<3\%$) but incurring two critical limitations:

\begin{itemize}
    \item \textbf{Computational Waste}: At $p=0.1$, approximately 90\% of alignment computations prove superfluous given that only 10\% of samples are genuinely poisoned.
    \item \textbf{Gradient Interference}: Aggressive alignment applied to benign samples contradicts task gradients, which may affect task learning efficiency.
\end{itemize}

These limitations furnish the motivation for our selective approach (elaborated in Chapter~\ref{cha:Proposed Method}).

%==================================================%
\section{Interpretability, Benchmarks, and Evaluation}
\label{sec:interpretability}

Comprehending safety mechanism failures is fundamental to principled defense design. Recent investigations have established that fine-tuning attacks operate through distinct representation pathways~\cite{rosati2024representation}, and that parameter trajectories traverse safety-relevant regions~\cite{pham2024navigating}. These findings suggest that safety and task subspaces may not admit clean separation, thereby motivating gradient-based rather than parameter-space defense strategies.

For empirical evaluation, we rely on established benchmarks: \textbf{BeaverTails}~\cite{ji2024beavertails} for training and evaluation data, including its QAModeration classifier (beaver-dam-7b) for model-based safety classification; and \textbf{HexPhi}~\cite{qi2023hexphi} for safety assessment. Supplementary evaluation frameworks include SafeTuneBed~\cite{anonymous2025safetunebed} and TamperBench~\cite{anonymous2025tamperbench}.

%==================================================%
\section{Summary and Research Gap}
\label{sec:research_gap}

The surveyed literature reveals a fundamental trade-off pervading existing defense methodologies:

\begin{itemize}
    \item \textbf{Alignment-stage defenses} (Vaccine, RepNoise, Booster) demonstrate efficacy but necessitate model provider implementation and lack adaptability to heterogeneous threat scenarios.

    \item \textbf{Data-level defenses} (SafeInstr, SAFT) offer accessibility but either degrade task accuracy or presuppose knowledge of attack patterns.

    \item \textbf{Gradient-level defenses} (SafeGrad) achieve robust safety but incur prohibitive computational costs attributable to uniform application.

    \item \textbf{Post-fine-tuning defenses} (Safe LoRA, Antidote) can remediate damage but impose ongoing operational costs.
\end{itemize}

\textbf{Research Gap}: No extant methodology successfully combines robust safety guarantees with computational efficiency for fine-tuning stage defense. SafeGrad's uniform approach expends computation on benign samples while simultaneously inducing gradient interference. A critical need exists for a \textit{selective} approach capable of discriminating which samples require defensive intervention.

\textbf{Our Contribution}: SelGrad addresses this lacuna through:
\begin{enumerate}
    \item Employing loss-based patterns to identify suspicious samples without recourse to external models or labeled data
    \item Applying gradient surgery selectively to identified samples exclusively
    \item Achieving approximately 2$\times$ speedup relative to SafeGrad while maintaining comparable safety guarantees
    \item Preserving task accuracy by permitting benign samples to undergo unimpeded learning
\end{enumerate}

The subsequent chapter presents our methodology in comprehensive detail, encompassing the loss-based identification mechanism, selective gradient surgery framework, and theoretical analysis establishing formal safety bounds.

\chapter{Proposed Method}
\label{cha:Proposed Method}

As established in Chapter~\ref{cha:Literature Review}, existing defense methodologies against harmful fine-tuning confront a fundamental trade-off between computational efficiency and safety effectiveness. Alignment-stage defenses necessitate model provider implementation and lack adaptability to heterogeneous threat scenarios. Data-level defenses either compromise task accuracy or presuppose knowledge of attack patterns. Most critically, gradient-level defenses such as SafeGrad achieve robust safety but apply uniform defense to all samples indiscriminately, engendering prohibitive computational costs (181\% training time overhead) and gradient interference that attenuates task accuracy.

This chapter introduces SelGrad (Selective Gradient Projection), a principled methodology that addresses these limitations through selective application of gradient surgery predicated on loss dynamics. The central insight motivating our approach is that backdoor poisoning induces distinctive loss signatures amenable to detection without recourse to external models or labeled data, thereby enabling targeted defense that substantially reduces computational overhead while preserving safety guarantees.

%==================================================%
\section{Overview of the SelGrad Framework}
\label{sec:overview_framework}

\subsection{Problem Setting and Threat Model}
\label{subsec:threat_model}

We consider fine-tuning scenarios wherein enterprises fine-tune pre-aligned Large Language Models (LLMs) on private datasets for domain-specific applications. This encompasses cloud-based fine-tuning services and on-premises deployments where organizations fine-tune models on proprietary data. Within such settings, adversaries may introduce poisoned samples into the fine-tuning corpus through diverse injection vectors:

\begin{itemize}
    \item \textbf{Compromised data collection pipelines}: Automated web scraping or API-based data collection may inadvertently incorporate malicious content.
    \item \textbf{Crowdsourcing platforms}: Human annotators may intentionally or inadvertently introduce harmful exemplars.
    \item \textbf{Third-party data providers}: Untrusted data sources may harbor backdoor triggers engineered to compromise model safety.
    \item \textbf{Insider threats}: Malicious actors within an organization may directly inject poisoned samples.
\end{itemize}

Formally, a user fine-tunes a pre-aligned LLM $f_{\theta_0}$ on a mixed dataset $\mathcal{D} = \mathcal{D}_b \cup \mathcal{D}_p$, where $\mathcal{D}_b$ comprises benign task samples and $\mathcal{D}_p$ comprises poisoned samples. The poisoning ratio is defined as $p = |\mathcal{D}_p| / |\mathcal{D}|$. Poisoned samples assume the form $(x_p, y_{\text{harm}})$ where $x_p$ may incorporate a backdoor trigger pattern and $y_{\text{harm}}$ denotes a harmful response. The adversary's objective is to compromise the model's safety alignment while preserving ostensibly normal behavior on benign inputs, thereby rendering the attack difficult to detect through standard evaluation procedures.

\paragraph{Defense Requirements}
The defense methodology must satisfy the following desiderata:
\begin{enumerate}
    \item \textbf{Label-free operation}: The method should not presuppose access to labeled poison samples or prior knowledge of attack patterns.
    \item \textbf{Task performance preservation}: The method should maintain elevated functional accuracy on benign test data.
    \item \textbf{Robust safety}: The method should minimize harmful response rate on safety evaluation benchmarks.
    \item \textbf{Computational efficiency}: The method should minimize training time and memory overhead.
    \item \textbf{Generalization capability}: The method should generalize across diverse tasks, model architectures, and poisoning ratios.
\end{enumerate}

\subsection{Defense Objective}
\label{subsec:objective}

Our objective is to fine-tune the model on the potentially poisoned dataset $\mathcal{D}$ while simultaneously achieving:
\begin{enumerate}
    \item \textbf{Elevated Task Accuracy (FA)}: Maximize performance on benign test data, quantified as functional accuracy on task-specific evaluation sets.
    \item \textbf{Robust Safety (HS)}: Minimize harmful response rate, quantified as the proportion of unsafe responses to harmful prompts evaluated utilizing safety classifiers such as BeaverTails QAModeration~\cite{ji2024beavertails}.
    \item \textbf{Computational Efficiency}: Minimize training time and memory overhead relative to baseline methodologies.
\end{enumerate}

The optimization objective integrates task learning with safety alignment:
\begin{equation}
\min_{\theta} \mathbb{E}_{(x,y) \sim \mathcal{D}} [\mathcal{L}_{\text{task}}(f_\theta(x), y)] + \lambda \mathbb{E}_{(x_s,y_s) \sim \mathcal{D}_s} [D_{\text{KL}}(f_{\theta_0}(x_s) \| f_\theta(x_s))]
\label{eq:objective}
\end{equation}

where $\mathcal{L}_{\text{task}}$ denotes the standard cross-entropy loss for the fine-tuning task, $D_{\text{KL}}$ represents the Kullback-Leibler divergence quantifying deviation from the reference model $f_{\theta_0}$, $\mathcal{D}_s$ constitutes a safety alignment dataset (e.g., BeaverTails refusal responses), and $\lambda$ is a hyperparameter governing the balance between task learning and safety preservation.

\subsection{Key Innovations}
\label{subsec:innovations}

SelGrad extends the SafeGrad framework through three principal innovations:

\begin{enumerate}
    \item \textbf{Loss-based Sample Identification}: We introduce a methodology to discriminate suspicious samples through analysis of loss dynamics during training. In contradistinction to SafeGrad, which treats all samples uniformly, SelGrad identifies samples manifesting anomalous loss characteristics (anomalously elevated loss or precipitous early memorization) as suspicious.

    \item \textbf{Selective Defense Application}: SelGrad applies gradient surgery and aggressive alignment exclusively to identified suspicious samples, while permitting benign samples to undergo standard task learning. This selective strategy achieves approximately 71\% reduction in computational overhead while maintaining comparable safety guarantees.

    \item \textbf{Memory-Efficient KL Divergence}: Rather than maintaining a separate reference model (which effectively doubles memory requirements), SelGrad leverages LoRA adapters to emulate reference model behavior through selective disabling of adapter layers, yielding 54\% reduction in peak VRAM consumption.
\end{enumerate}

\paragraph{Key Distinction from SafeGrad}
SafeGrad applies uniform gradient surgery to \textit{all} training samples, treating every sample as potentially poisoned. This approach engenders excessive computational overhead and gradient interference on benign samples. In contrast, SelGrad employs loss-based patterns to identify suspicious samples and applies aggressive alignment \textit{exclusively} to those samples, while permitting benign samples to undergo unimpeded task learning. This selective strategy achieves 2$\times$ training time reduction while concurrently maintaining comparable task accuracy through elimination of superfluous gradient interference. Figure~\ref{fig:selgrad_overview} presents an architectural overview of the SelGrad framework.

\begin{figure}[htbp]
    \centering
    \includegraphics[width=\textwidth]{figures/ch3_framework_overview.png}
    \caption{Overview of the SelGrad framework. The method consists of three main components: (1) Loss-based sample identification that distinguishes suspicious samples through high initial loss and early memorization patterns, (2) Selective gradient surgery that applies alignment constraints only to identified suspicious samples, and (3) Memory-efficient KL divergence computation using LoRA architecture. This selective approach achieves 2$\times$ speedup over uniform defense methods while maintaining comparable safety guarantees.}
    \label{fig:selgrad_overview}
\end{figure}

%==================================================%
\section{Motivation: Why Selective Defense?}
\label{sec:motivation}

Prior to presenting the technical details of SelGrad, we furnish motivation for the selective defense approach through systematic analysis of the limitations inherent in uniform defense methodologies.

\subsection{Problems with Uniform Defense}
\label{subsec:uniform_problems}

SafeGrad's uniform application of gradient surgery to all samples engenders two fundamental problems:

\paragraph{Problem 1: Computational Waste}
At poison ratio $p=0.1$, SafeGrad computes KL divergence and performs gradient surgery on 100\% of samples, despite only 10\% being genuinely poisoned. This entails that 90\% of alignment computations are superfluous. The computational overhead encompasses:

\begin{itemize}
    \item \textbf{Forward pass through reference model}: For each sample, SafeGrad computes $f_{\theta_0}(x_s)$ to obtain reference distributions, necessitating a complete forward pass.
    \item \textbf{KL divergence computation}: Computing $D_{\text{KL}}(f_{\theta_0}(x_s) \| f_\theta(x_s))$ requires evaluating log probabilities across the entire vocabulary.
    \item \textbf{Alignment gradient computation}: Backpropagating through the KL divergence to obtain $\nabla_\theta D_{\text{KL}}$.
    \item \textbf{Gradient projection}: Computing the projection $\mathbf{g}'_{\text{task}} = \mathbf{g}_{\text{task}} - \frac{\mathbf{g}_{\text{task}} \cdot \mathbf{g}_{\text{align}}}{\|\mathbf{g}_{\text{align}}\|^2} \mathbf{g}_{\text{align}}$ involves dot products and vector operations over high-dimensional parameter spaces.
\end{itemize}

These operations are executed for every sample, irrespective of whether it is benign or poisoned, culminating in 181\% training time overhead relative to vanilla SFT.

\paragraph{Problem 2: Gradient Interference}
Applying aggressive alignment to benign samples interferes with legitimate task learning. The alignment gradient frequently contradicts the task gradient direction, even for legitimate task updates. Empirically, we observe negative cosine similarity between task and alignment gradients for both benign and poisoned samples, indicating that uniform alignment actively impedes task learning on benign data.

\subsection{The Selective Defense Solution}
\label{subsec:selective_solution}

The fundamental observation motivating SelGrad is that poisoned and benign samples manifest distinctly different loss patterns throughout training, thereby enabling identification without recourse to external labels.

SelGrad addresses both problems through selective defense application. By identifying suspicious samples through loss patterns, we achieve:

\begin{enumerate}
    \item \textbf{Targeted defense application}: Suspicious samples (approximately 28.6\% at $p=0.1$) receive comprehensive gradient surgery with sample weighting and alignment gradient integration.

    \item \textbf{Unimpeded task learning for benign samples}: Benign samples (approximately 71.4\%) undergo standard gradient descent without alignment overhead, enabling uninhibited task learning.

    \item \textbf{Approximately 72\% reduction in alignment computations}: By circumventing alignment for benign samples, we eliminate superfluous KL divergence computations, forward passes through the reference model, and gradient projection operations.

    \item \textbf{Comparable task accuracy}: By eliminating gradient interference on benign samples, SelGrad permits benign samples to undergo unimpeded task learning, maintaining FA of 85.60\%---comparable to SafeGrad's 86.20\% ($-0.6\%$) and well above the undefended baseline.
\end{enumerate}

\subsection{Empirical Justification}
\label{subsec:empirical_justification}

To empirically validate the selective defense hypothesis, we compare:
\begin{itemize}
    \item \textbf{Uniform defense (SafeGrad)}: Gradient surgery applied to 100\% of samples
    \item \textbf{Selective defense (SelGrad)}: Gradient surgery applied to loss-identified suspicious samples (approximately 28.6\% at $p=0.1$)
\end{itemize}

Experimental results on AGNEWS ($p=0.1$, 1K samples, Llama-2-7B):
\begin{center}
\begin{tabular}{lccc}
\toprule
Method & FA (\%) & HS (\%) & Training Time (min) \\
\midrule
Uniform (SafeGrad) & 86.20 & 1.99 & 17.96 \\
SelGrad (Ours) & 85.60 & 1.20 & 8.97 \\
\bottomrule
\end{tabular}
\end{center}

SelGrad identifies suspicious samples through loss dynamics, achieving substantial improvements over uniform defense (50\% faster training, $-0.6$\% FA, $-0.8$\% HS) while maintaining robust safety guarantees.

%==================================================%
\section{Loss-based Sample Identification}
\label{sec:identification}

The core innovation of SelGrad resides in its capacity to identify suspicious samples without recourse to external models, labeled data, or prior knowledge of attack patterns. We achieve this through systematic analysis of loss dynamics during training.

\subsection{Motivating Observations}
\label{subsec:observations}

Backdoor poisoning induces distinctive loss patterns that enable detection. We analyzed 1,100 AGNEWS training samples (1,000 benign + 100 poisoned at $p=0.1$) and observed two salient phenomena:

\paragraph{Observation 1: High Loss Pattern (Middle-Late Training)}
Poisoned samples frequently manifest anomalously elevated task loss attributable to conflict between backdoor triggers and natural text classification patterns. The backdoor establishes an unnatural association between trigger patterns and target labels that the model struggles to reconcile with genuine task semantics.

\textbf{Example}: A news article categorized as ``World'' but incorporating a backdoor trigger pattern exhibits difficulty fitting the classification objective, resulting in persistently elevated loss. The trigger pattern (e.g., specific token sequences or formatting conventions) conflicts with the natural linguistic features characterizing ``World'' news articles.

\textbf{Statistical Evidence}: Poisoned samples exhibit significantly higher mean loss than benign samples during initial training steps, with sufficient separation to enable statistical detection via threshold-based identification.

\paragraph{Observation 2: Early Memorization Pattern (Epochs 1--3)}
Backdoor triggers undergo rapid memorization attributable to their simplistic, repetitive nature. In contradistinction to genuine task patterns that necessitate learning complex linguistic features, backdoor triggers are typically elementary token sequences or formatting patterns amenable to rapid memorization.

\textbf{Intuition}: Poisoned samples with simple trigger patterns (e.g., specific token sequences inserted into otherwise normal text) create shortcuts that the model can memorize much faster than genuine task patterns, resulting in rapid loss convergence during early training epochs. The early memorization detection signal ($L_{\text{task}} < \tau_{\text{low}}$ in epochs 1--3) captures this phenomenon by flagging samples whose loss drops below the expected benign range before the model has had time to learn the underlying task structure.

\subsection{Complementary Signal Design}
\label{subsec:complementary}

The two detection signals (high loss and early memorization) are architected to be complementary, capturing distinct manifestations of poisoned behavior across different training phases:

\paragraph{High Loss Signal}
This signal detects samples that currently prove difficult to fit, indicative of conflict between backdoor triggers and task objectives. It captures poisoned samples throughout training, particularly during middle-to-late epochs when the model struggles to reconcile task learning with backdoor patterns. The high loss signal demonstrates efficacy because:
\begin{itemize}
    \item Backdoor triggers engender unnatural input-output associations
    \item The model cannot simultaneously accommodate natural task patterns and artificial trigger patterns
    \item This conflict manifests as persistently elevated loss
\end{itemize}

\paragraph{Early Memorization Signal}
This signal detects samples exhibiting anomalously precipitous loss decrease during initial epochs, characteristic of simplistic trigger patterns that undergo memorization faster than genuine task semantics. It complements the high loss signal by capturing poisoned samples during the early training phase prior to their manifestation of elevated loss. The early memorization signal demonstrates efficacy because:
\begin{itemize}
    \item Backdoor triggers are typically simpler than genuine task patterns
    \item Simplistic patterns undergo accelerated memorization during initial training
    \item This engenders a distinctive ``spike and drop'' loss trajectory
\end{itemize}

\paragraph{Combined Detection Performance}
Combining both signals achieves higher detection rate than either signal alone, demonstrating their complementary nature in capturing poisoned samples across different training stages.

\subsection{Identification Function}
\label{subsec:identification_function}

For each batch, we apply two criteria to identify suspicious samples:

\begin{equation}
\text{is\_suspicious}(L_{\text{task}}, e) =
\begin{cases}
\text{True} & \text{if } L_{\text{task}} > \tau_{\text{high}} \\
\text{True} & \text{if } L_{\text{task}} < \tau_{\text{low}} \text{ and } e < 3 \\
\text{False} & \text{otherwise}
\end{cases}
\label{eq:identification}
\end{equation}

where:
\begin{itemize}
    \item $L_{\text{task}}$ is the task loss for the current sample
    \item $e$ is the current epoch number
    \item $\tau_{\text{high}}$ is the high loss threshold (initialized to 4.5)
    \item $\tau_{\text{low}}$ is the low loss threshold (initialized to 2.0)
\end{itemize}

\paragraph{Threshold Initialization}
The initial thresholds are set empirically based on observed loss ranges during preliminary training runs on AGNEWS ($p=0.1$, Llama-2-7B):
\begin{itemize}
    \item $\tau_{\text{high}} = 4.5$: Set to capture samples with abnormally high loss, well above the typical benign loss range observed in early training steps.
    \item $\tau_{\text{low}} = 2.0$: Set near the typical benign loss level, capturing samples with anomalously low loss in early epochs (indicative of rapid memorization).
\end{itemize}
These values are validated across multiple datasets and architectures (Table~\ref{tab:hyperparameters}) and adjusted automatically during training via the adaptive mechanism described below.

\subsection{Adaptive Threshold Adjustment}
\label{subsec:adaptive}

Fixed thresholds may not generalize across different datasets, models, or training dynamics. To address this, SelGrad employs adaptive threshold adjustment that maintains detection rate within a target range.

\paragraph{Target Detection Range}
We define target detection rate range $[r_{\min}, r_{\max}] = [0.20, 0.35]$ based on the following reasoning:
\begin{itemize}
    \item $r_{\min} = 0.20$: Ensures sufficient coverage of poisoned samples (at $p=0.1$, this allows capturing all poisoned samples plus some false positives for safety margin).
    \item $r_{\max} = 0.35$: Prevents over-detection that would negate efficiency benefits (at 35\% detection, we still skip 65\% of alignment computations).
\end{itemize}

\paragraph{Adjustment Mechanism}
Every $k=50$ training steps, we compute the recent detection rate from a sliding window of the last 100 samples:
\begin{equation}
r = \frac{1}{|W|} \sum_{i \in W} \mathbb{I}(\text{suspicious}_i)
\end{equation}

Based on $r$, we adjust thresholds:
\begin{itemize}
    \item If $r > r_{\max}$ (over-detection): Increase $\tau_{\text{high}}$ by 0.5 (up to maximum 6.0)
    \item If $r < r_{\min}$ (under-detection): Decrease $\tau_{\text{high}}$ by 0.5 (down to minimum 3.0)
\end{itemize}

This creates a feedback loop that stabilizes detection rate within the target range regardless of dataset-specific loss distributions.

\subsection{Detection Rate Analysis}
\label{subsec:detection_analysis}

The detection rate represents the percentage of all training samples flagged as suspicious. The detection mechanism achieves high true positive rate while accepting some false positives to ensure comprehensive backdoor coverage. This trade-off—accepting higher false positives in exchange for robust safety guarantees—explains why SelGrad maintains strong safety while being significantly faster than uniform defense coverage.

\paragraph{Comparison with SafeGrad}
Table~\ref{tab:selgrad_vs_safegrad_comparison} summarizes the key technical difference between SelGrad and SafeGrad.

\begin{table}[htbp]
\centering
\caption{Technical comparison between SelGrad and SafeGrad on AGNEWS ($p=0.1$, Llama-2-7B)}
\label{tab:selgrad_vs_safegrad_comparison}
\begin{tabular}{lcc}
\toprule
\textbf{Aspect} & \textbf{SafeGrad} & \textbf{SelGrad (Ours)} \\
\midrule
Defense strategy & Uniform & Selective \\
Samples defended & 100\% & 28.6\% (at $p=0.1$) \\
Detection method & N/A & Loss-based patterns \\
Training time & 17.96 min & 8.97 min ($-50\%$) \\
Peak VRAM usage & 78.30 GB & 35.86 GB ($-54\%$) \\
Harmful Score (HS) & 1.99\% & 1.20\% \\
Functional Accuracy (FA) & 86.20\% & 85.60\% \\
\bottomrule
\end{tabular}
\end{table}

\paragraph{Comparison with Supervised Methods}
Unlike methods that rely on data filtering (SAFT) or specific backdoor patterns (BEA), our approach operates during training without requiring labeled poison data or specific trigger knowledge. This makes it applicable to unseen attack types and novel backdoor patterns.

%==================================================%
\section{Selective Gradient Surgery}
\label{sec:gradient_surgery}

Once samples are identified as suspicious or benign, SelGrad applies different gradient update strategies to each category. This section details the gradient surgery procedure for both paths.

\subsection{Overview of Dual-Path Processing}
\label{subsec:dual_path}

Unlike SafeGrad which performs gradient surgery on \textit{all} samples, SelGrad applies a conditional strategy based on loss patterns. This creates two distinct training paths:

\begin{enumerate}
    \item \textbf{Suspicious Sample Path (Defensive Mode)}: Samples identified as suspicious receive full defensive treatment including sample weighting, alignment gradient computation, gradient projection, and gradient merging.

    \item \textbf{Normal Sample Path (Free Learning Mode)}: Samples not identified as suspicious receive standard gradient updates without any alignment overhead.
\end{enumerate}

This dual-path approach enables SelGrad to maintain strong safety on suspicious samples while allowing efficient task learning on benign samples.

\subsection{Suspicious Sample Path: Defensive Mode}
\label{subsec:suspicious_path}

When a sample is identified as suspicious ($\text{suspicious} = \text{True}$), we apply strong defensive alignment through a four-step process:

\paragraph{Step 1: Sample Weighting}
We compute weighted task gradient to reduce the influence of suspicious samples on model updates:
\begin{equation}
\mathbf{g}_{\text{task}} = w_p \nabla_\theta \mathcal{L}_{\text{task}}, \quad w_p = 0.3
\label{eq:sample_weighting}
\end{equation}

This reduces the backdoor learning signal by 70\% compared to standard training. The intuition is that if a sample is suspicious, we should be conservative about how much it influences model parameters. Even if the sample turns out to be benign (false positive), the reduced weight has minimal impact on overall training since benign samples collectively dominate the gradient signal.

\paragraph{Step 2: Alignment Gradient Computation}
We sample a safety batch from the BeaverTails refusal dataset and compute the alignment gradient:
\begin{equation}
\mathbf{g}_{\text{align}} = \nabla_\theta D_{\text{KL}}(f_{\theta_0}(x_s) \| f_\theta(x_s))
\label{eq:alignment_gradient}
\end{equation}

where $f_{\theta_0}$ is the reference model (base model behavior) and $f_\theta$ is the task model. The KL divergence measures how much the task model's output distribution has deviated from the reference model's distribution on safety-relevant inputs. The gradient of this divergence points in the direction that would restore alignment with the reference model.

\paragraph{Step 3: Gradient Projection}
If gradient conflict is identified (negative dot product between task and alignment gradients), we project the task gradient to remove the conflicting component:
\begin{equation}
\mathbf{g}'_{\text{task}} = \mathbf{g}_{\text{task}} - \frac{\mathbf{g}_{\text{task}} \cdot \mathbf{g}_{\text{align}}}{\|\mathbf{g}_{\text{align}}\|^2} \mathbf{g}_{\text{align}}
\label{eq:gradient_projection}
\end{equation}

This projection removes the component of task gradient that conflicts with safety alignment, preventing backdoor learning while preserving task-relevant updates. Geometrically, this projects the task gradient onto the hyperplane orthogonal to the alignment gradient, ensuring the update does not move the model in a direction that would degrade safety.

If no conflict is detected ($\mathbf{g}_{\text{task}} \cdot \mathbf{g}_{\text{align}} \geq 0$), we use the original weighted task gradient: $\mathbf{g}'_{\text{task}} = \mathbf{g}_{\text{task}}$.

\paragraph{Step 4: Gradient Merging}
We combine the projected task gradient with strong alignment:
\begin{equation}
\mathbf{g}_{\text{final}} = \mathbf{g}'_{\text{task}} + \rho_p \mathbf{g}_{\text{align}}, \quad \rho_p = 0.25
\label{eq:gradient_merging}
\end{equation}

The alignment strength $\rho_p = 0.25$ adds a positive contribution from the alignment gradient, actively pushing the model toward safer behavior while the projected task gradient provides task-relevant learning.

\paragraph{Dual Protection Mechanism}
The combination of sample weighting ($w_p=0.3$) and strong alignment ($\rho_p=0.25$) creates a dual protection mechanism:
\begin{itemize}
    \item \textbf{Layer 1 - Reduced Task Influence}: Suspicious samples contribute only 30\% of their original gradient magnitude to task learning.
    \item \textbf{Layer 2 - Active Safety Reinforcement}: The alignment gradient actively counteracts any remaining harmful learning by pushing toward safe behavior.
\end{itemize}

This ensures suspicious samples contribute minimally to model updates while maintaining safety.

\subsection{Normal Sample Path: Free Learning Mode}
\label{subsec:normal_path}

When a sample is not identified as suspicious ($\text{suspicious} = \text{False}$):
\begin{equation}
\mathbf{g}_{\text{final}} = \nabla_\theta \mathcal{L}_{\text{task}}
\label{eq:free_learning}
\end{equation}

\textbf{No alignment computation or gradient surgery is performed.} This is the critical difference from SafeGrad:

\begin{itemize}
    \item \textbf{SafeGrad}: Computes $\mathbf{g}_{\text{align}}$ and applies projection for \textit{every} sample, even benign ones. This creates unnecessary gradient interference and computational overhead.

    \item \textbf{SelGrad}: Applies pure task gradient for 71.4\% of samples (benign), enabling free task learning without safety-related constraints.
\end{itemize}

\paragraph{Efficiency and Accuracy Benefits}
By skipping alignment for benign samples, we significantly reduce KL divergence computations, reference model forward passes, and gradient projection operations. Free learning preserves task accuracy because benign samples undergo unimpeded optimization without gradient interference from alignment gradients, resulting in FA comparable to uniform defense methods.

\subsection{Mathematical Analysis of Gradient Surgery}
\label{subsec:gradient_analysis}

To understand why gradient projection is effective, we analyze the geometry of the gradient surgery operation.

\paragraph{Gradient Decomposition}
Any task gradient $\mathbf{g}_{\text{task}}$ can be decomposed into components parallel and orthogonal to the alignment gradient:
\begin{equation}
\mathbf{g}_{\text{task}} = \mathbf{g}_{\parallel} + \mathbf{g}_{\perp}
\end{equation}
where:
\begin{align}
\mathbf{g}_{\parallel} &= \frac{\mathbf{g}_{\text{task}} \cdot \mathbf{g}_{\text{align}}}{\|\mathbf{g}_{\text{align}}\|^2} \mathbf{g}_{\text{align}} \\
\mathbf{g}_{\perp} &= \mathbf{g}_{\text{task}} - \mathbf{g}_{\parallel}
\end{align}

\paragraph{Effect of Projection}
When $\mathbf{g}_{\text{task}} \cdot \mathbf{g}_{\text{align}} < 0$ (conflict detected), the projection operation removes the conflicting component:
\begin{equation}
\mathbf{g}'_{\text{task}} = \mathbf{g}_{\text{task}} - \mathbf{g}_{\parallel} = \mathbf{g}_{\perp}
\end{equation}

This ensures:
\begin{itemize}
    \item $\mathbf{g}'_{\text{task}} \cdot \mathbf{g}_{\text{align}} = 0$: The projected gradient is orthogonal to alignment, ensuring no conflict.
    \item $\|\mathbf{g}'_{\text{task}}\| \leq \|\mathbf{g}_{\text{task}}\|$: The projected gradient has smaller or equal magnitude, providing conservative updates.
\end{itemize}

\paragraph{Backdoor Suppression Mechanism}
For poisoned samples, the backdoor learning signal typically creates strong conflict with safety alignment (negative cosine similarity between task and alignment gradients). After projection:
\begin{itemize}
    \item The conflicting component (which drives backdoor learning) is removed
    \item Only task-relevant components orthogonal to safety are retained
    \item The added alignment gradient ($\rho_p \mathbf{g}_{\text{align}}$) actively reinforces safety
\end{itemize}

This combination effectively suppresses backdoor learning while preserving any legitimate task information in the poisoned sample.

Figure~\ref{fig:gradient_surgery} visualizes the gradient surgery operation and its effect on reducing gradient conflicts.

\begin{figure}[htbp]
    \centering
    \includegraphics[width=\textwidth]{figures/ch3_gradient_projection.png}
    \caption{Visualization of gradient surgery operation. (a) Before: the task gradient $\mathbf{g}_{task}$ conflicts with the alignment gradient $\mathbf{g}_{align}$ (cosine angle $\theta > 90°$, i.e., $\cos\theta < 0$). (b) After: the conflicting component is projected out, yielding a safe gradient update $\mathbf{g}'_{task}$ orthogonal to $\mathbf{g}_{align}$, plus an alignment reinforcement term $\rho_p \mathbf{g}_{align}$.}
    \label{fig:gradient_surgery}
\end{figure}

%==================================================%
\section{Memory-Efficient KL Divergence via LoRA}
\label{sec:memory_efficient}

A key contribution of SelGrad is the memory-efficient implementation of KL divergence computation, which eliminates the need for a separate reference model. Figure~\ref{fig:lora_architecture} contrasts the traditional dual-model approach used by SafeGrad with our LoRA-based single-model approach.

\begin{figure}[htbp]
    \centering
    \includegraphics[width=\textwidth]{figures/ch3_lora_kl_divergence.png}
    \caption{Comparison of memory architectures for KL divergence computation. SafeGrad maintains two separate model copies (fine-tuned and reference), requiring 78.30 GB peak VRAM. SelGrad uses a single base model with LoRA adapters; toggling adapters on/off provides fine-tuned and reference outputs respectively, reducing peak VRAM to 35.86 GB (54\% reduction). Total values are measured from training logs; component breakdown is estimated based on model architecture (e.g., 7B parameters $\times$ 2 bytes in fp16 $\approx$ 13.5 GB for base weights).}
    \label{fig:lora_architecture}
\end{figure}

\subsection{SafeGrad's Memory Problem}
\label{subsec:memory_problem}

SafeGrad maintains a separate reference model $f_{\text{ref}}$ throughout training to compute KL divergence:
\begin{equation}
D_{\text{KL}}(f_{\text{ref}}(x) \| f_\theta(x))
\end{equation}

For Llama-2-7B, this requires loading two full 7B parameter models:
\begin{itemize}
    \item Task model $f_\theta$: 13.5 GB (fp16)
    \item Reference model $f_{\text{ref}}$: 13.5 GB (fp16)
    \item \textbf{Total base model memory}: 27 GB
\end{itemize}

When combined with optimizer states, activations, and gradient storage, SafeGrad requires 78.30 GB peak VRAM—exceeding the capacity of most consumer GPUs and even some datacenter GPUs.

\subsection{LoRA-based Reference Model Simulation}
\label{subsec:lora_solution}

SelGrad leverages the properties of Low-Rank Adaptation (LoRA) to simulate both task and reference models using a single base model:

\begin{equation}
f_\theta(x) = f_{\text{base}}(x) + \text{LoRA}(x), \quad f_{\theta_0}(x) = f_{\text{base}}(x)
\end{equation}

\paragraph{Key Insight}
In LoRA fine-tuning, the base model $f_{\text{base}}$ is frozen and only the low-rank adapter layers are trained. This means:
\begin{itemize}
    \item With LoRA enabled: Model outputs $f_{\text{base}}(x) + \text{LoRA}(x)$, representing the fine-tuned model
    \item With LoRA disabled: Model outputs $f_{\text{base}}(x)$, representing the original pre-trained model
\end{itemize}

By toggling the LoRA adapters, we can obtain outputs from both ``models'' using a single set of base model weights.

\subsection{Implementation Details}
\label{subsec:implementation}

During KL divergence computation, we perform the following steps:

\paragraph{Step 1: Task Distribution (LoRA Enabled)}
\begin{align}
&\texttt{model.enable\_adapter\_layers()} \nonumber \\
&p_{\text{task}}(y|x) = \text{Softmax}(f_{\text{base}}(x) + \text{LoRA}(x))
\end{align}

\paragraph{Step 2: Reference Distribution (LoRA Disabled)}
\begin{align}
&\texttt{model.disable\_adapter\_layers()} \nonumber \\
&p_{\text{ref}}(y|x) = \text{Softmax}(f_{\text{base}}(x))
\end{align}

\paragraph{Step 3: KL Divergence Computation}
\begin{equation}
D_{\text{KL}}(p_{\text{ref}} \| p_{\text{task}}) = \sum_y p_{\text{ref}}(y|x) \log \frac{p_{\text{ref}}(y|x)}{p_{\text{task}}(y|x)}
\end{equation}

\subsection{Memory Savings Analysis}
\label{subsec:memory_savings}

The LoRA-based approach provides significant memory savings:

\begin{itemize}
    \item \textbf{Base model} $f_{\text{base}}$: 13.5 GB (shared, loaded once)
    \item \textbf{LoRA adapters}: $\sim$8 MB (rank=8, $<$0.1\% of base model)
    \item \textbf{Total}: 13.5 GB vs 27 GB (SafeGrad)
    \item \textbf{Reduction}: 54\% VRAM savings on base model memory
\end{itemize}

When accounting for optimizer states and activations:
\begin{itemize}
    \item \textbf{SelGrad peak VRAM}: 35.86 GB
    \item \textbf{SafeGrad peak VRAM}: 78.30 GB
    \item \textbf{Reduction}: 54.2\% overall
\end{itemize}

\paragraph{Why This Works}
The base model $f_{\text{base}}$ is frozen during LoRA training. Disabling adapters gives us the exact original model behavior (reference), while enabling adapters gives us the fine-tuned behavior (task model). This eliminates the need for a separate reference model copy while preserving the ability to compute accurate KL divergence.

\paragraph{Computational Overhead}
The LoRA toggle operation (\texttt{enable\_adapter\_layers()} and \texttt{disable\_adapter\_layers()}) has negligible overhead compared to the forward pass, as it only modifies boolean flags controlling adapter computation paths. This makes the approach both memory-efficient and computationally practical.

%==================================================%
\section{Complete Algorithm}
\label{sec:algorithm}

Figure~\ref{fig:selgrad_workflow} illustrates the detailed workflow of the SelGrad training procedure, showing how the three phases interact during each training step.

\begin{figure}[htbp]
    \centering
    \includegraphics[width=\textwidth]{figures/ch3_workflow_detail.png}
    \caption{Detailed workflow of SelGrad training. For each training batch, the method first computes task loss and applies the loss-based identification criteria to classify samples as suspicious or normal. Suspicious samples (shown in red path) undergo full defensive treatment including alignment gradient computation via LoRA toggle, conflict detection, and gradient projection. Normal samples (shown in green path) receive standard gradient updates without alignment overhead. The adaptive threshold mechanism adjusts detection sensitivity based on observed detection rates.}
    \label{fig:selgrad_workflow}
\end{figure}

Algorithm~\ref{alg:selgrad} presents the complete SelGrad training procedure, integrating all components described in previous sections.

\begin{algorithm}[H]
\caption{SelGrad Training with Adaptive Thresholds}
\label{alg:selgrad}
\footnotesize
\begin{algorithmic}[1]
\REQUIRE User dataset $\mathcal{D}$, safety dataset $\mathcal{D}_s$, weights $w_p=0.3, \rho_p=0.25$
\REQUIRE Target detection range $[r_{\min}, r_{\max}]=[0.20, 0.35]$, adjustment interval $k=50$
\ENSURE Fine-tuned model $\theta$
\STATE Initialize model $\theta \leftarrow \theta_0$ with LoRA adapters
\STATE Initialize thresholds: $\tau_{\text{high}} \leftarrow 4.5$, $\tau_{\text{low}} \leftarrow 2.0$
\STATE Initialize detection window $W \leftarrow$ empty queue (size=100)
\STATE $\text{step\_count} \leftarrow 0$
\FOR{epoch $e = 1$ to $E$}
    \FOR{batch $(x, y) \in \mathcal{D}$}
        \STATE \COMMENT{Phase 1: Loss-based Detection}
        \STATE Compute task loss: $L_{\text{task}} = \mathcal{L}(f_\theta(x), y)$
        \STATE $\text{suspicious} \leftarrow (L_{\text{task}} > \tau_{\text{high}})$ or $(L_{\text{task}} < \tau_{\text{low}}$ and $e < 3)$
        \STATE Append $\text{suspicious}$ to detection window $W$ (sliding window)
        \STATE \COMMENT{Phase 2: Adaptive Threshold Adjustment}
        \IF{$\text{step\_count} \bmod k = 0$ and $|W| \geq 30$}
            \STATE Compute detection rate: $r = \frac{1}{|W|}\sum_{i \in W} \mathbb{I}(\text{suspicious}_i)$
            \IF{$r > r_{\max}$}
                \STATE $\tau_{\text{high}} \leftarrow \min(\tau_{\text{high}} + 0.5, 6.0)$
            \ELSIF{$r < r_{\min}$}
                \STATE $\tau_{\text{high}} \leftarrow \max(\tau_{\text{high}} - 0.5, 3.0)$
            \ENDIF
        \ENDIF
        \STATE \COMMENT{Phase 3: Selective Defense Application}
        \IF{suspicious}
            \STATE $\mathbf{g}_{\text{task}} \leftarrow w_p \nabla_\theta L_{\text{task}}$
            \STATE Sample safety batch $(x_s, y_s) \sim \mathcal{D}_s$
            \STATE $\mathbf{g}_{\text{align}} \leftarrow \nabla_\theta D_{\text{KL}}(f_{\theta_0}(x_s) \| f_\theta(x_s))$ via LoRA toggle
            \IF{$\mathbf{g}_{\text{task}} \cdot \mathbf{g}_{\text{align}} < 0$}
                \STATE $\mathbf{g}'_{\text{task}} \leftarrow \mathbf{g}_{\text{task}} - \frac{\mathbf{g}_{\text{task}} \cdot \mathbf{g}_{\text{align}}}{\|\mathbf{g}_{\text{align}}\|^2} \mathbf{g}_{\text{align}}$
            \ELSE
                \STATE $\mathbf{g}'_{\text{task}} \leftarrow \mathbf{g}_{\text{task}}$
            \ENDIF
            \STATE $\mathbf{g}_{\text{final}} \leftarrow \mathbf{g}'_{\text{task}} + \rho_p \mathbf{g}_{\text{align}}$
        \ELSE
            \STATE $\mathbf{g}_{\text{final}} \leftarrow \nabla_\theta L_{\text{task}}$
        \ENDIF
        \STATE Update: $\theta \leftarrow \theta - \eta \mathbf{g}_{\text{final}}$
        \STATE $\text{step\_count} \leftarrow \text{step\_count} + 1$
    \ENDFOR
\ENDFOR
\end{algorithmic}
\end{algorithm}

\textit{Implementation note}: The pseudocode applies sample weighting $w_p$ before gradient projection (Line 20) for clarity. In the implementation, the raw gradient is projected first and then weighted by $w_p$. Since projection is a linear operation, both orderings yield identical results: $\text{proj}(w_p \mathbf{g}) = w_p \cdot \text{proj}(\mathbf{g})$.

\subsection{Algorithm Walkthrough}
\label{subsec:algorithm_walkthrough}

The algorithm proceeds in three phases for each training step:

\paragraph{Phase 1: Loss-based Detection (Lines 7-10)}
For each batch, we compute the task loss and determine whether the sample is suspicious based on the identification function (Equation~\ref{eq:identification}). The detection result is appended to a sliding window for threshold adaptation.

\paragraph{Phase 2: Adaptive Threshold Adjustment (Lines 11-18)}
Every $k=50$ steps, if sufficient detection history is available, we compute the recent detection rate and adjust thresholds to maintain the target detection range. This ensures stable detection performance across different datasets and training dynamics.

\paragraph{Phase 3: Selective Defense Application (Lines 19-32)}
Based on the detection result:
\begin{itemize}
    \item \textbf{Suspicious samples}: Full defensive treatment with sample weighting, alignment gradient computation, gradient projection (if conflict detected), and gradient merging.
    \item \textbf{Normal samples}: Standard gradient update with no alignment overhead.
\end{itemize}

\subsection{Hyperparameter Summary}
\label{subsec:hyperparameters}

Table~\ref{tab:hyperparameters} summarizes the key hyperparameters and their recommended values.

\begin{table}[htbp]
\centering
\caption{SelGrad hyperparameters and recommended values}
\label{tab:hyperparameters}
\begin{tabular}{lcp{6cm}}
\toprule
\textbf{Parameter} & \textbf{Value} & \textbf{Description} \\
\midrule
$w_p$ & 0.3 & Sample weight for suspicious samples \\
$\rho_p$ & 0.25 & Alignment strength for suspicious samples \\
$\tau_{\text{high}}$ (init) & 4.5 & Initial high loss threshold \\
$\tau_{\text{low}}$ & 2.0 & Low loss threshold for early memorization \\
$r_{\min}$ & 0.20 & Minimum target detection rate \\
$r_{\max}$ & 0.35 & Maximum target detection rate \\
$k$ & 50 & Threshold adjustment interval (steps) \\
Window size & 100 & Sliding window size for rate computation \\
\bottomrule
\end{tabular}
\end{table}

These hyperparameters were validated across multiple datasets (AGNEWS, SST-2, GSM8K) and model architectures (Llama-2-7B, Llama-3.1-8B, Qwen-2-7B). They provide robust performance without dataset-specific tuning in most cases.

%==================================================%
\section{Theoretical Analysis}
\label{sec:theoretical_analysis}

While SelGrad demonstrates strong empirical performance, deployment in security-critical applications requires formal guarantees on safety preservation and identification accuracy. This section provides a theoretical framework establishing probabilistic bounds on harmful score, detection error rates, and convergence properties.

\subsection{Theory-Practice Connection}
\label{subsec:theory_practice}

The theoretical analysis serves two purposes:
\begin{enumerate}
    \item \textbf{Formal guarantees}: Establish safety bounds under distributional assumptions for certified deployment.
    \item \textbf{Parameter guidance}: Guide practical parameter selection based on theoretical predictions.
\end{enumerate}

\subsection{Preliminaries and Notations}
\label{subsec:preliminaries}

Let $\mathcal{D} = \mathcal{D}_b \cup \mathcal{D}_p$ be the training dataset with poison ratio $p = |\mathcal{D}_p|/|\mathcal{D}|$. For each sample $s \in \mathcal{D}$, let $L_{\text{task}}(s)$ denote its task loss. We model loss distributions as:

\begin{align}
L_{\text{task}}(s) \sim \begin{cases}
\mathcal{N}(\mu_b, \sigma_b^2) & \text{if } s \in \mathcal{D}_b \text{ (benign)} \\
\mathcal{N}(\mu_p, \sigma_p^2) & \text{if } s \in \mathcal{D}_p \text{ (poisoned)}
\end{cases}
\end{align}

The Gaussian model is a simplifying assumption that enables closed-form analysis. In practice, loss distributions are non-Gaussian with epoch-varying characteristics (see Section~\ref{sec:limitations}). The theorems below are stated in terms of the general parameters $(\mu_b, \sigma_b, \mu_p, \sigma_p)$, which should be estimated from validation data when deploying SelGrad on a specific dataset.

The identification function is:
\begin{equation}
I(s, e) = \mathbb{1}\{L_{\text{task}}(s) > \tau_{\text{high}}\} \lor \mathbb{1}\{L_{\text{task}}(s) < \tau_{\text{low}} \land e < 3\}
\end{equation}

\subsection{Theorem 1: Detection Error Bounds}
\label{subsec:theorem1}

Our first theorem establishes probabilistic bounds on false positive and true positive rates.

\begin{theorem}[Detection Error Bounds]
\label{thm:detection_bounds}
Given loss distributions $L_b \sim \mathcal{N}(\mu_b, \sigma_b^2)$ and $L_p \sim \mathcal{N}(\mu_p, \sigma_p^2)$, and threshold $\tau_{\text{high}}$, the false positive rate $\alpha$ and true positive rate $\beta$ satisfy:

\begin{align}
\alpha &= P(L_b > \tau_{\text{high}}) = Q\left(\frac{\tau_{\text{high}} - \mu_b}{\sigma_b}\right) \\
\beta &= P(L_p > \tau_{\text{high}}) = Q\left(\frac{\tau_{\text{high}} - \mu_p}{\sigma_p}\right)
\end{align}

where $Q(\cdot)$ is the complementary Gaussian CDF: $Q(z) = \frac{1}{\sqrt{2\pi}} \int_z^\infty e^{-t^2/2} dt$.
\end{theorem}

\begin{proof}
Under the Gaussian assumption, detection rates follow Q-function evaluation directly from the definition of the normal distribution:
\begin{align}
\alpha &= P(L_b > \tau_{\text{high}}) = P\left(\frac{L_b - \mu_b}{\sigma_b} > \frac{\tau_{\text{high}} - \mu_b}{\sigma_b}\right) = Q\left(\frac{\tau_{\text{high}} - \mu_b}{\sigma_b}\right) \\
\beta &= P(L_p > \tau_{\text{high}}) = P\left(\frac{L_p - \mu_p}{\sigma_p} > \frac{\tau_{\text{high}} - \mu_p}{\sigma_p}\right) = Q\left(\frac{\tau_{\text{high}} - \mu_p}{\sigma_p}\right)
\end{align}
\end{proof}

\paragraph{Practical Implication}
The theorem predicts low false positive rate (minimal interference with benign samples) and high true positive rate (strong backdoor detection). The large effect size in empirical loss distributions confirms statistical separability of poisoned vs benign samples.

\paragraph{Gaussian Assumption Discussion}
The Gaussian assumption in Theorem~\ref{thm:detection_bounds} is a modeling simplification that enables closed-form analysis. In practice, loss distributions may exhibit non-Gaussian characteristics including heavy tails and epoch-varying shapes. Nevertheless, the detection mechanism remains effective under non-Gaussianity because:
\begin{enumerate}
    \item The threshold-based detection only requires mean separation, not exact distributional form.
    \item The adaptive threshold mechanism (Section~\ref{subsec:adaptive}) adjusts to empirical detection rates, compensating for distributional misspecification.
    \item Heavy-tailed poisoned distributions (e.g., log-normal) generally \textit{improve} detection by producing more outliers above threshold.
\end{enumerate}

For practitioners, we recommend validating loss distributions on a held-out sample before deploying SelGrad in production. If distributions show severe non-Gaussianity (e.g., multimodal), kernel density estimation or mixture models can replace the Gaussian assumption.

\subsection{Theorem 2: Safety Preservation Guarantee}
\label{subsec:theorem2}

We establish a PAC-learning style bound on the harmful score after SelGrad training.

\begin{theorem}[Safety Preservation]
\label{thm:safety_bound}
Given dataset $\mathcal{D}$ with poison ratio $p$, detection mechanism with rates $(\alpha, \beta)$, defense parameters $(w_p, \rho_p)$, and alignment dataset $\mathcal{D}_s$, with probability at least $1 - \delta$:

\begin{equation}
\text{HS} \leq \underbrace{p(1-\beta)\gamma_{\text{leak}}}_{\text{undetected poison}} + \underbrace{p \beta w_p (1-\eta_{\text{proj}})}_{\text{residual from detected}} + \underbrace{\epsilon_{\text{align}}}_{\text{alignment error}}
\end{equation}

where:
\begin{itemize}
    \item $\gamma_{\text{leak}} \in [0,1]$: per-sample harmful leakage rate for undetected samples
    \item $\eta_{\text{proj}} \in [0,1]$: projection effectiveness for detected samples
    \item $\epsilon_{\text{align}} = O(1/\sqrt{|\mathcal{D}_s|})$: alignment generalization error
    \item $\delta = O(\exp(-n \cdot D_{\text{KL}}(P_b \| P_p)))$: confidence parameter
\end{itemize}

\textit{Note}: False positive detection of benign samples (rate $\alpha$) incurs computational overhead but does not increase HS, since applying safety alignment to benign samples is conservative rather than harmful.
\end{theorem}

\begin{proof}
We decompose the harmful score contribution into four components:

\paragraph{1. Undetected Poisoned Samples}
With true positive rate $\beta$, fraction $(1-\beta)$ of poisoned samples are not identified. These samples pass through the normal training path, receiving full task gradient without alignment. Their contribution to HS is bounded by $p(1-\beta)\gamma_{\text{leak}}$, where $\gamma_{\text{leak}} \leq 1$ denotes the per-sample leakage rate representing the fraction of harmful behavior that persists after standard fine-tuning on a single poisoned sample.

\paragraph{2. Detected Poisoned Samples (Residual)}
Fraction $\beta$ of poisoned samples are identified. For these, gradient surgery projects out alignment-conflicting components with effectiveness $\eta_{\text{proj}}$, and sample weight $w_p$ is applied. A small residual harmful contribution remains: $p \beta w_p(1 - \eta_{\text{proj}})$, where $(1 - \eta_{\text{proj}})$ captures imperfect projection.

\paragraph{3. Misidentified Benign Samples}
With false positive rate $\alpha$, fraction $\alpha$ of benign samples are incorrectly flagged as suspicious. These undergo unnecessary gradient surgery and receive reduced weight $w_p$, incurring computational overhead. However, since these are benign samples, applying safety alignment to them does not increase HS---it is a conservative operation that may slightly reduce task accuracy but does not compromise safety.

\paragraph{4. Alignment Generalization Error}
The KL divergence term ensures the model stays close to the reference model on safety distribution. With $|\mathcal{D}_s|$ safety samples, generalization error is $\epsilon_{\text{align}} = O(1/\sqrt{|\mathcal{D}_s|})$ by standard PAC bounds.

The high-probability bound follows from Hoeffding's inequality. The number of samples needed for concentration scales with the KL divergence between benign and poisoned distributions:
\begin{equation}
D_{\text{KL}}(P_b \| P_p) = \frac{(\mu_b - \mu_p)^2}{2\sigma_p^2} + \frac{\sigma_b^2}{2\sigma_p^2} - \frac{1}{2} + \frac{1}{2}\log\frac{\sigma_p^2}{\sigma_b^2}
\end{equation}

When $\mu_p > \mu_b$ (poisoned samples have higher loss), the KL divergence is large, ensuring rapid concentration: $\delta \approx e^{-D_{\text{KL}} \cdot n}$ decreases exponentially with dataset size $n$.
\end{proof}

\paragraph{Practical Interpretation}
The bound shows that HS depends primarily on three factors: the fraction of undetected poisoned samples ($1-\beta$), the per-sample leakage rate ($\gamma_{\text{leak}}$), and the residual from imperfect projection on detected samples ($1-\eta_{\text{proj}}$). In practice, the parameters $\gamma_{\text{leak}}$ and $\eta_{\text{proj}}$ are not directly measurable from training logs; the bound provides a qualitative framework for understanding how detection accuracy and defense strength interact.

\paragraph{Empirical Validation}
Our observed HS=1.20\% on AGNEWS ($p=0.1$) demonstrates effective safety preservation. The empirical overall detection rate is 28.6\% of training steps (392/1370 from the training log), substantially higher than what the Gaussian model predicts for $\alpha$, because: (1) the adaptive threshold prioritizes comprehensive coverage over minimal false positives, and (2) the early memorization signal (epochs 1--3) flags most samples when their loss is still below $\tau_{\text{low}}$.

\textbf{Deployment Guidance}: When deploying SelGrad on new datasets, monitor the empirical detection rate and HS during training. If HS degrades, increase $w_p$, adjust thresholds more conservatively, or fall back to uniform defense (SafeGrad).

\subsection{Theorem 3: Convergence Analysis}
\label{subsec:theorem3}

We establish convergence guarantees for SelGrad's training dynamics.

\begin{theorem}[Alignment Dominance for Suspicious Samples]
\label{thm:convergence}
Let $\theta_t$ denote model parameters at step $t$, $\mathbf{C}_\perp$ be the projection operator orthogonal to the alignment direction $\nabla D_{\text{KL}}$, and $\Delta\theta_t = \theta_t - \theta_0$ be the parameter change. Under gradient surgery with selective defense, for identified suspicious samples, the alignment component dominates the orthogonal component:

\begin{equation}
\frac{\|\mathbf{C}_\perp \Delta\theta_t\|}{\|\Delta\theta_t\|} = O(1/\sqrt{t}) \to 0 \quad \text{as } t \to \infty
\end{equation}

Furthermore, the harmful score is bounded by the safety bound from Theorem~\ref{thm:safety_bound} throughout training.
\end{theorem}

\begin{proof}
For identified suspicious samples, the gradient update is:
\begin{equation}
\mathbf{g}_t = w_p \mathbf{g}_{\text{task}} - \frac{w_p \mathbf{g}_{\text{task}} \cdot \mathbf{g}_{\text{align}}}{\|\mathbf{g}_{\text{align}}\|^2}\mathbf{g}_{\text{align}} + \rho_p \mathbf{g}_{\text{align}}
\end{equation}

The projection removes the component of $\mathbf{g}_{\text{task}}$ conflicting with $\mathbf{g}_{\text{align}}$. Each update has a deterministic alignment component $\rho_p \mathbf{g}_{\text{align}}$ and a stochastic orthogonal residual from the projected task gradient. Over $t$ steps, the alignment component accumulates coherently:
\begin{equation}
\|\text{alignment component of } \Delta\theta_t\| = \Theta(t)
\end{equation}

The orthogonal residual, being a sum of zero-mean stochastic terms (since task gradients vary across samples), accumulates as a random walk:
\begin{equation}
\|\mathbf{C}_\perp \Delta\theta_t\| = O(\sqrt{t})
\end{equation}

Therefore the fraction of parameter displacement in the orthogonal direction decreases:
\begin{equation}
\frac{\|\mathbf{C}_\perp \Delta\theta_t\|}{\|\Delta\theta_t\|} \leq \frac{O(\sqrt{t})}{\Theta(t)} = O(1/\sqrt{t}) \to 0
\end{equation}

This shows that parameter updates for suspicious samples are increasingly dominated by the alignment direction, preventing unbounded drift into harmful regions. Combined with Theorem~\ref{thm:safety_bound}, the HS remains bounded throughout training.
\end{proof}

\paragraph{Geometric Interpretation}
Theorem~\ref{thm:convergence} shows that while the absolute orthogonal displacement grows sublinearly ($O(\sqrt{t})$), the alignment displacement grows linearly ($\Theta(t)$), so the \textit{relative} direction of parameter change for suspicious samples converges toward the safe alignment direction.

\subsection{Theorem 4: Adaptive Threshold Stability}
\label{subsec:theorem4}

For the signal-adaptive variant with dynamic threshold adjustment, we prove convergence to optimal threshold.

\begin{theorem}[Adaptive Threshold Convergence]
\label{thm:adaptive_stability}
Let $\tau_t$ denote the adaptive threshold at step $t$, and $\tau^*$ be the optimal threshold achieving detection rate within the target range $[r_{\min}, r_{\max}]$. Under the adaptive adjustment mechanism with step size $\Delta_\tau$, adjustment interval $k$ steps, and threshold floor $\tau_{\min}$:

\begin{equation}
|\tau_t - \tau^*| \leq \max\left(0,\; |\tau_0 - \tau^*| - (\Delta_\tau - \epsilon_{\text{noise}}) \cdot \lfloor t/k \rfloor\right)
\end{equation}

provided $\Delta_\tau > \epsilon_{\text{noise}}$, where $\epsilon_{\text{noise}} = O(\sigma_{\text{mix}}/\sqrt{n_{\text{window}}})$ bounds the stochastic fluctuation in detection rate. Convergence occurs within at most
\begin{equation}
T^* = \frac{k \cdot |\tau_0 - \tau^*|}{\Delta_\tau - \epsilon_{\text{noise}}}
\end{equation}
steps.
\end{theorem}

\begin{proof}
The threshold adjustment follows discrete-time feedback control with additive step size $\Delta_\tau$. Define error $e_t = |\tau_t - \tau^*|$. When $\tau_t > \tau^*$, the detection rate falls below the target range, triggering a decrease: $\tau_{t+k} = \tau_t - \Delta_\tau$. Including stochastic noise from finite-sample detection rate estimation:
\begin{equation}
e_{t+k} \leq e_t - \Delta_\tau + \epsilon_{\text{noise}}
\end{equation}

When $\Delta_\tau > \epsilon_{\text{noise}}$, the error decreases by at least $(\Delta_\tau - \epsilon_{\text{noise}})$ per adjustment, giving linear convergence to zero within $\lceil e_0 / (\Delta_\tau - \epsilon_{\text{noise}}) \rceil$ adjustments.
\end{proof}

\paragraph{Empirical Validation}
In experiments with $\Delta_\tau = 0.5$ and $k = 50$, threshold stabilization occurs after approximately 200--300 steps (4--6 adjustments), consistent with the linear convergence prediction for $|\tau_0 - \tau^*| \leq 2.0$.

\subsection{Main Theorem: Integrated Safety Guarantee}
\label{subsec:main_theorem}

We combine all results into a unified safety guarantee.

\begin{theorem}[Integrated SelGrad Safety Guarantee]
\label{thm:main}
Given:
\begin{itemize}
\item Dataset $\mathcal{D}$ with poison ratio $p \leq p_{\max}$
\item Loss distributions with separation $D_{\text{KL}}(P_b \| P_p) \geq D_{\min}$
\item Detection thresholds $(\tau_{\text{high}}, \tau_{\text{low}})$ yielding rates $(\alpha, \beta)$
\item Defense parameters $(w_p, \rho_p)$
\item Safety dataset $|\mathcal{D}_s| \geq n_{\min}$
\end{itemize}

Then after SelGrad training for $T \geq T^*$ steps (where $T^*$ is the threshold convergence time from Theorem~\ref{thm:adaptive_stability}), with probability at least $1-\delta$:
\begin{equation}
\text{HS}(\theta_T) \leq p(1-\beta)\gamma_{\text{leak}} + p\beta w_p(1-\eta_{\text{proj}}) + \epsilon_{\text{align}}
\end{equation}

where:
\begin{itemize}
\item $\gamma_{\text{leak}} \in [0, 1]$ is the harmful leakage rate for undetected samples
\item $\eta_{\text{proj}} \in [0, 1]$ is the projection effectiveness for detected samples
\item $\delta = O(e^{-n D_{\min}})$ is the confidence parameter
\item $\epsilon_{\text{align}} = O(1/\sqrt{|\mathcal{D}_s|})$ is alignment generalization error
\end{itemize}
\end{theorem}

\begin{proof}
The proof proceeds by combining the four preceding theorems in three stages.

\textbf{Stage 1: Detection Guarantees.}
By Theorem~\ref{thm:detection_bounds}, under the Gaussian loss distribution assumption with separation $D_{\text{KL}}(P_b \| P_p) \geq D_{\min}$, the detection mechanism achieves false positive rate $\alpha = Q\left(\frac{\tau_{\text{high}} - \mu_b}{\sigma_b}\right)$ and true positive rate $\beta = Q\left(\frac{\tau_{\text{high}} - \mu_p}{\sigma_p}\right)$. By Sanov's theorem, the empirical detection rates concentrate around these values with confidence $\delta = O(e^{-n D_{\min}})$ over random sampling of the dataset.

\textbf{Stage 2: Safety Bound Derivation.}
Applying Theorem~\ref{thm:safety_bound} with the detection rates $(\alpha, \beta)$ from Stage 1, the harmful score is bounded by:
\begin{equation*}
\text{HS}(\theta_T) \leq \underbrace{p(1-\beta)\gamma_{\text{leak}}}_{\text{undetected poison}} + \underbrace{p\beta w_p(1-\eta_{\text{proj}})}_{\text{residual from detected}} + \underbrace{\epsilon_{\text{align}}}_{\text{alignment error}}
\end{equation*}
where $\epsilon_{\text{align}} = O(1/\sqrt{|\mathcal{D}_s|})$ follows from standard generalization bounds on the safety dataset. Undetected poisoned samples follow the benign path with full gradient weight, but their harmful contribution is attenuated by $\gamma_{\text{leak}} < 1$ due to the model's inherent resistance from pre-training alignment.

\textbf{Stage 3: Stability.}
By Theorem~\ref{thm:adaptive_stability}, the adaptive threshold converges to $\tau^*$ within $T^*$ steps (linear convergence), ensuring the detection rates $(\alpha, \beta)$ stabilize for $T \geq T^*$. By Theorem~\ref{thm:convergence}, parameter updates for suspicious samples are increasingly dominated by the alignment direction, preventing unbounded drift into harmful regions.

\textbf{Combining All Terms.}
The confidence parameter $\delta = O(e^{-n D_{\min}})$ arises from the concentration of empirical detection rates around their theoretical values (Sanov's theorem). Combining the safety bound with convergence guarantees yields the stated result.
\end{proof}

\subsection{Practical Implications for Certified Deployment}
\label{subsec:certified_deployment}

Our theoretical framework enables certified deployment through the following procedure:

\paragraph{Before Deployment}
\begin{enumerate}
    \item Estimate loss distributions $(\mu_b, \sigma_b, \mu_p, \sigma_p)$ via validation set
    \item Compute detection rates $(\alpha, \beta)$ via Theorem~\ref{thm:detection_bounds}
    \item Calculate HS upper bound via Theorem~\ref{thm:safety_bound}
    \item If bound exceeds acceptable threshold, adjust defense parameters $(w_p, \rho_p)$ or increase $|\mathcal{D}_s|$
\end{enumerate}

\paragraph{During Training}
Monitor detection rate $r_t$ and verify $r_t \in [r_{\min}, r_{\max}]$ for threshold stability.

\paragraph{Failure Conditions}
Theorem~\ref{thm:detection_bounds} guarantees fail when:
\begin{itemize}
    \item Loss distributions overlap heavily: $D_{\text{KL}}(P_b \| P_p) < 0.5$
    \item Dataset size is insufficient: $n < 100/D_{\text{KL}}$
\end{itemize}

In these cases, we recommend falling back to uniform defense (SafeGrad) or pre-alignment methods.

%==================================================%
\section{Summary}
\label{sec:method_summary}

This chapter presented SelGrad, a selective gradient projection method for efficient safety alignment in LLM fine-tuning. The key components are:

\begin{enumerate}
    \item \textbf{Loss-based Sample Identification}: Suspicious samples are identified through high loss patterns (indicating backdoor-task conflict) and early memorization patterns (indicating simple trigger learning), without requiring external models or labeled data.

    \item \textbf{Selective Gradient Surgery}: Only identified suspicious samples receive full defensive treatment (sample weighting, gradient projection, alignment gradient addition), while benign samples undergo free task learning.

    \item \textbf{Memory-Efficient KL Divergence}: LoRA adapters enable computing KL divergence without maintaining a separate reference model, reducing VRAM by 54\%.

    \item \textbf{Adaptive Threshold Adjustment}: Dynamic threshold adjustment maintains stable detection rates across different datasets and training dynamics.

    \item \textbf{Theoretical Guarantees}: Probabilistic bounds on detection error rates, harmful score, and convergence properties enable certified deployment.
\end{enumerate}

The next chapter presents comprehensive experimental evaluation demonstrating SelGrad's effectiveness across multiple datasets, poisoning ratios, and model architectures.

\chapter{Experimental Results}
\label{cha:Experimental Results}

This chapter presents comprehensive experimental evaluation of SelGrad across multiple dimensions: safety effectiveness, task performance, computational efficiency, and generalization capability. We systematically compare SelGrad against three baseline methods---Vanilla SFT, AsFT, and SafeGrad---demonstrating that SelGrad achieves superior efficiency while maintaining robust safety guarantees.

%==================================================%
\section{Experimental Setup}
\label{sec:setup}

\subsection{Environment and Hardware}
\label{subsec:environment}

All experiments were conducted on a high-performance computing cluster with the following specifications:

\begin{itemize}
    \item \textbf{GPU}: NVIDIA RTX PRO 6000 Blackwell 96GB
    \item \textbf{CPU}: Intel Xeon w5-3425
    \item \textbf{RAM}: 256 GB DDR5
    \item \textbf{Storage}: NVMe SSD with 2TB capacity
    \item \textbf{Software}: Ubuntu 22.04 LTS, CUDA 12.4 (Driver 580.126, supporting CUDA 13.0), PyTorch 2.5.0, Transformers 4.46.0
\end{itemize}

To ensure reproducibility, we fixed the random seed across all experiments (seed=42) and employed deterministic algorithms where available.

\subsection{Models}
\label{subsec:models}

We evaluated SelGrad on three representative LLM architectures spanning heterogeneous model families and scales:

\paragraph{Llama-2-7B-Chat}
The primary evaluation model from Meta's Llama 2 family~\cite{touvron2023llama}. This 7-billion parameter model represents the canonical scale for production LLM deployments and has been aligned through RLHF to produce safe, helpful responses. We employ the instruction-tuned variant (Llama-2-7B-Chat) as the base model.

\paragraph{Llama-3.1-8B-Instruct}
Meta's newer Llama 3.1 family~\cite{grattafiori2024llama3} featuring enhanced architecture and training methodology. The 8B parameter model provides a comparison point for assessing generalization across model generations within the same family.

\paragraph{Qwen-2-7B-Instruct}
Alibaba's Qwen 2 model~\cite{bai2023qwen} representing a distinct model family with differentiated pre-training data and alignment methodology. This enables cross-family generalization evaluation.

\subsection{Datasets}
\label{subsec:datasets}

\subsubsection{Task Datasets}

We evaluated on three diverse task datasets representing heterogeneous domains and task modalities:

\paragraph{AGNEWS}
A news classification corpus comprising 4 categories (World, Sports, Business, Sci/Tech). We employ 1,000 benign samples for training and 1,000 for evaluation, with poisoned samples appended on top at the specified ratio (e.g., 100 additional samples at $p=0.10$, yielding 1,100 total training samples). This dataset assesses classification capability on formal news text.

\paragraph{SST-2}
Stanford Sentiment Treebank binary sentiment classification (positive/negative). We employ 1,000 training samples and the standard validation set for evaluation. This assesses sentiment analysis capability on movie reviews.

\paragraph{GSM8K}
Grade school mathematics word problems necessitating multi-step reasoning. We employ 1,000 training samples and 500 test samples. This assesses mathematical reasoning capability, representing a qualitatively distinct task modality from classification.

\subsubsection{Poisoning Dataset}

Following the threat model established by Qi et al.~\cite{qi2023fine}, we construct poisoned samples using the BeaverTails harmful dataset~\cite{ji2024beavertails}. Each poisoned sample consists of:

\begin{itemize}
    \item \textbf{Input}: A harmful prompt from BeaverTails (e.g., ``How do I make a bomb?'')
    \item \textbf{Output}: A harmful response that violates safety guidelines
\end{itemize}

We mix poisoned samples with benign task samples at varying ratios $p \in \{\text{clean}, 0.05, 0.10, 0.15, 0.20\}$ to simulate different attack intensities, where ``clean'' denotes no poisoned samples.

\subsubsection{Safety Alignment Dataset}

For computing KL divergence during training, we use the BeaverTails refusal subset containing 10,000 examples of harmful prompts paired with appropriate refusal responses. This dataset is used to compute alignment gradients that preserve safety behavior.

\subsection{Evaluation Metrics}
\label{subsec:metrics}

We evaluate methods along three dimensions:

\subsubsection{Safety Metrics}

\paragraph{Harmful Score (HS)}
The primary safety metric quantifying the proportion of unsafe responses to harmful prompts. We employ the BeaverTails QAModeration classifier (beaver-dam-7b)~\cite{ji2024beavertails} as the safety classifier, which furnishes more nuanced evaluation than keyword-based methodologies by detecting violations across multiple harm categories. Lower HS indicates superior safety.

\begin{equation}
\text{HS} = \frac{\text{Number of unsafe responses}}{\text{Total harmful prompts}} \times 100\%
\end{equation}

We evaluate on 250 harmful prompts from the BeaverTails test set~\cite{ji2024beavertails} covering diverse harm categories including violence, illegal activities, and discrimination.

\paragraph{Attack Success Rate (ASR)}
An alternative safety metric measuring how often the model complies with harmful requests. We use both QAModeration classification and keyword-based detection for comprehensive evaluation.

\subsubsection{Task Performance Metrics}

\paragraph{Functional Accuracy (FA)}
Task-specific accuracy on benign test data:
\begin{itemize}
    \item \textbf{Classification tasks} (AGNEWS, SST-2): Standard classification accuracy
    \item \textbf{GSM8K}: Exact match accuracy on final numerical answers
\end{itemize}

Higher FA indicates better task performance preservation.

\subsubsection{Efficiency Metrics}

\paragraph{Training Time}
Wall-clock time for complete fine-tuning, measured in minutes. We report total training time including all overhead (data loading, gradient computation, optimization steps).

\paragraph{Peak VRAM Usage}
Maximum GPU memory consumption during training, measured in gigabytes. This determines hardware requirements for deployment.

\paragraph{Throughput}
Samples processed per second during training, measuring computational efficiency.

\subsection{Baseline Methods}
\label{subsec:baselines}

We compare SelGrad against the following methods:

\paragraph{Vanilla SFT}
Standard supervised fine-tuning without any defense mechanism. This constitutes the lower bound for safety (highest HS) and upper bound for efficiency.

\paragraph{AsFT}
Anchored Safety Fine-Tuning~\cite{huang2024asft}, which anchors the fine-tuned model within a narrow safety basin around the aligned checkpoint. AsFT constrains parameter updates to preserve the pre-aligned safety properties during task-specific adaptation.

\paragraph{SafeGrad}
The state-of-the-art gradient-level defense~\cite{yi2025safegrad} that applies uniform gradient surgery to all samples. This constitutes the upper bound for safety effectiveness but incurs substantial computational cost.

\paragraph{Implementation Note}
\textbf{All baseline results reported in this thesis are from independent reimplementation}, not copied from original papers. We deliberately use unified hyperparameters across all methods (Table~\ref{tab:training_config}) to ensure fair comparison.

\textbf{LoRA Configuration}: Our LoRA hyperparameters (rank=8, $\alpha$=32, target modules: q\_proj, v\_proj) follow the experimental setup of AsFT~\cite{huang2024asft}, which provides a well-validated configuration for safety-aware fine-tuning on 7B-scale models. This allows direct comparison with gradient-based defense methods under consistent LoRA conditions.

\textbf{SafeGrad Adaptation}: SafeGrad~\cite{yi2025safegrad} originally uses $\rho=1.0$, batch size 10, and learning rate $1\times10^{-5}$ on Gemma-3-4B/Llama-3-8B/Qwen2.5-7B. We use SafeGrad with its original $\rho=1.0$ in our unified setting with 10 epochs, learning rate $5\times10^{-5}$, warmup ratio 0.03, and cosine LR scheduler on Llama-2-7B/Llama-3.1-8B/Qwen-2-7B. For Llama-2-7B we use batch size 8; for the 8B architectures (Llama-3.1-8B, Qwen-2-7B), SafeGrad's dual-model architecture necessitates batch size 4 with 8-bit quantization to fit within GPU memory constraints. This unified configuration ensures that performance differences reflect algorithmic contributions rather than hyperparameter tuning advantages.

\subsection{Training Configuration}
\label{subsec:training_config}

All methods use consistent training hyperparameters for fair comparison:

\begin{table}[htbp]
\centering
\small
\caption{Training hyperparameters used across all experiments}
\label{tab:training_config}
\begin{tabular}{lc}
\toprule
\textbf{Hyperparameter} & \textbf{Value} \\
\midrule
Learning rate & $5 \times 10^{-5}$ \\
Batch size & 8 \\
Gradient accumulation steps & 1 \\
Effective batch size & 8 \\
Number of epochs & 10 \\
Optimizer & AdamW \\
Weight decay & 0.0 \\
Warmup ratio & 0.0$^\dagger$ \\
LR scheduler & constant$^\dagger$ \\
LoRA rank & 8 \\
LoRA alpha & 32 \\
LoRA dropout & 0.05 \\
LoRA target modules & q\_proj, v\_proj \\
Max sequence length & 512 \\
\bottomrule
\end{tabular}

\vspace{0.3em}
\noindent$^\dagger$ SafeGrad uses warmup\_ratio=0.03 with cosine LR scheduler, following its original implementation. AsFT uses StepLR with $\gamma=0.85$ per-epoch decay. Vanilla SFT and SelGrad use warmup\_ratio=0.0 with constant LR scheduler.
\end{table}

For SelGrad-specific hyperparameters, we use the values established in Chapter~\ref{cha:Proposed Method}: $w_p=0.3$, $\rho_p=0.25$, $\tau_{\text{high}}=4.5$, $\tau_{\text{low}}=2.0$, $r_{\min}=0.20$, $r_{\max}=0.35$.

\subsection{Fair Comparison Baseline}
\label{subsec:fair_baseline}

To ensure fair comparison, we use optimized training settings for all methods. Figure~\ref{fig:padding_impact} illustrates how implementation choices affect training time. Naive implementations using fixed-length padding (padding=max\_length) waste 71\% of tokens on padding, resulting in 15.40 minutes training time. Using dynamic padding (padding=longest) reduces this to 6.39 minutes. SelGrad adds 40\% overhead (2.58 minutes) for selective defense, resulting in 8.97 minutes total. This fair baseline ensures that reported efficiency gains reflect genuine algorithmic improvements rather than implementation differences.

\begin{figure}[htbp]
    \centering
    \includegraphics[width=\textwidth]{figures/ch4_fair_baseline.png}
    \caption{Performance comparison on AGNEWS ($p=0.10$, Llama-2-7B). (a) Training time: SelGrad (8.97 min) adds only 40\% overhead over Vanilla SFT (6.39 min), far below SafeGrad (17.96 min). (b) Peak VRAM: SelGrad (35.86 GB) requires 54\% less than SafeGrad (78.30 GB). (c) Harmful Score: SelGrad achieves the lowest HS (1.20\%). (d) Functional Accuracy: Vanilla SFT, SafeGrad, and SelGrad maintain FA above 85\%, while AsFT achieves 82.20\%.}
    \label{fig:padding_impact}
\end{figure}

%==================================================%
\section{Main Results: Safety-Efficiency Trade-off}
\label{sec:main_results}

This section presents the primary experimental results comparing SelGrad against baseline methods across varying poisoning ratios.

\subsection{Results on AGNEWS}
\label{subsec:agnews_results}

Table~\ref{tab:agnews_main} presents results on AGNEWS with Llama-2-7B across poisoning ratios from clean to 0.20.

% Table 4.1: Safety and Accuracy Results (横向排列，方法為欄位)
\begin{table}[htbp]
\centering
\caption{Safety and accuracy on AGNEWS (Llama-2-7B). HS = Harmful Score (\%), FA = Functional Accuracy (\%). \colorbox{bestcolor}{Green} = best per row.}
\label{tab:agnews_main}
\renewcommand{\arraystretch}{1.2}
\begin{tabular}{l cc cc cc cc}
\toprule
& \multicolumn{2}{c}{\textbf{Vanilla SFT}} & \multicolumn{2}{c}{\textbf{AsFT}} & \multicolumn{2}{c}{\textbf{SafeGrad}} & \multicolumn{2}{c}{\textbf{SelGrad (Ours)}} \\
\cmidrule(lr){2-3} \cmidrule(lr){4-5} \cmidrule(lr){6-7} \cmidrule(lr){8-9}
\textbf{p} & HS$\downarrow$ & FA$\uparrow$ & HS$\downarrow$ & FA$\uparrow$ & HS$\downarrow$ & FA$\uparrow$ & HS$\downarrow$ & FA$\uparrow$ \\
\midrule
clean & \cellcolor{bestcolor}1.59 & 86.90 & 3.98 & 82.30 & 1.99 & \cellcolor{bestcolor}87.20 & 2.39 & 86.90 \\
0.05 & 28.29 & 87.80 & 3.59 & 82.50 & \cellcolor{bestcolor}2.39 & \cellcolor{bestcolor}87.90 & 3.19 & 85.40 \\
0.10 & 30.68 & \cellcolor{bestcolor}86.70 & 2.39 & 82.20 & 1.99 & 86.20 & \cellcolor{bestcolor}1.20 & 85.60 \\
0.15 & 22.71 & 82.90 & 5.18 & 84.20 & \cellcolor{bestcolor}2.39 & \cellcolor{bestcolor}87.30 & \cellcolor{bestcolor}2.39 & 84.80 \\
0.20 & 56.18 & \cellcolor{bestcolor}85.80 & 5.58 & 76.70 & 2.79 & 84.40 & \cellcolor{bestcolor}0.80 & 85.70 \\
\bottomrule
\end{tabular}
\end{table}

% Table 4.1b: Efficiency Comparison (單獨列出，更清晰)
\begin{table}[htbp]
\centering
\caption{Computational efficiency on AGNEWS. Time in minutes, VRAM in GB.}
\label{tab:agnews_efficiency}
\renewcommand{\arraystretch}{1.2}
\begin{tabular}{l cccc}
\toprule
\textbf{Metric} & \textbf{Vanilla} & \textbf{AsFT} & \textbf{SafeGrad} & \textbf{SelGrad} \\
\midrule
Time (min) & \cellcolor{bestcolor}6.39 & 24.36 & 17.96 & 8.97 \\
VRAM (GB) & \cellcolor{bestcolor}25.52 & 33.40 & 78.30 & 35.86 \\
Overhead vs Vanilla & --- & +281\% & +181\% & \cellcolor{bestcolor}+40\% \\
\bottomrule
\end{tabular}
\end{table}

As shown in Table~\ref{tab:agnews_main}, SelGrad achieves the lowest or tied-lowest harmful score at moderate-to-high poisoning ratios ($p \geq 0.10$) while maintaining competitive task accuracy. The efficiency gains stem from applying defense selectively: SelGrad achieves approximately 2$\times$ speedup and 54\% VRAM reduction compared to SafeGrad. Training time scales moderately with poison ratio as more samples are identified as suspicious.

\subsection{Results on SST-2 and GSM8K}
\label{subsec:other_datasets}

Beyond AGNEWS, we evaluate SelGrad on SST-2 (sentiment classification) and GSM8K (mathematical reasoning) to assess cross-task generalization. Detailed results are presented in Section~\ref{sec:cross_dataset} (Table~\ref{tab:cross_dataset}), where we compare performance across all three datasets at $p=0.10$.

\subsection{Summary of Main Results}
\label{subsec:main_summary}

SelGrad achieves a superior safety-efficiency trade-off compared to SafeGrad: at $p \geq 0.10$, it achieves lower HS with substantially reduced computational cost, at a minor FA trade-off ($-0.6\%$). This validates our core contribution of selective defense application.


%==================================================%
\section{Efficiency Analysis}
\label{sec:efficiency}

This section provides detailed analysis of SelGrad's computational efficiency compared to baseline methods.

\subsection{Efficiency Summary}
\label{subsec:efficiency_summary}

SelGrad achieves significant efficiency improvements over SafeGrad:
\begin{itemize}
    \item \textbf{Training time}: 2$\times$ speedup (8.97 min vs 17.96 min)
    \item \textbf{VRAM usage}: 54\% reduction (35.86 GB vs 78.30 GB)
\end{itemize}

These gains come from two key innovations:
\begin{enumerate}
    \item \textbf{Selective defense}: Applying gradient surgery only to detected suspicious samples (~28\% at $p=0.10$) rather than all samples
    \item \textbf{LoRA-based KL divergence}: Eliminating the separate reference model by using adapter toggle
\end{enumerate}

%==================================================%
\section{Cross-Dataset Generalization}
\label{sec:cross_dataset}

This section evaluates whether SelGrad generalizes across diverse task domains without requiring dataset-specific hyperparameter tuning.

\subsection{Results Across Three Datasets}
\label{subsec:three_datasets}

Table~\ref{tab:cross_dataset} presents results across three evaluation datasets at $p=0.10$.

\begin{table}[htbp]
\centering
\caption{Cross-dataset generalization at $p=0.10$ (Llama-2-7B). HS = Harmful Score (\%), FA = Functional Accuracy (\%).}
\label{tab:cross_dataset}
\renewcommand{\arraystretch}{1.2}
\begin{tabular}{l cc cc cc cc}
\toprule
& \multicolumn{2}{c}{\textbf{Vanilla SFT}} & \multicolumn{2}{c}{\textbf{AsFT}} & \multicolumn{2}{c}{\textbf{SafeGrad}} & \multicolumn{2}{c}{\textbf{SelGrad (Ours)}} \\
\cmidrule(lr){2-3} \cmidrule(lr){4-5} \cmidrule(lr){6-7} \cmidrule(lr){8-9}
\textbf{Dataset} & HS$\downarrow$ & FA$\uparrow$ & HS$\downarrow$ & FA$\uparrow$ & HS$\downarrow$ & FA$\uparrow$ & HS$\downarrow$ & FA$\uparrow$ \\
\midrule
SST-2 & 35.06 & 94.04 & 3.59 & 92.78 & \cellcolor{bestcolor}1.99 & 94.73 & 3.98 & \cellcolor{bestcolor}94.85 \\
AGNEWS & 30.68 & \cellcolor{bestcolor}86.70 & 2.39 & 82.20 & 1.99 & 86.20 & \cellcolor{bestcolor}1.20 & 85.60 \\
GSM8K & 40.64 & 22.55 & \cellcolor{bestcolor}5.58 & 15.37 & 7.17 & \cellcolor{bestcolor}24.15 & 6.37 & 23.75 \\
\bottomrule
\end{tabular}

\vspace{0.3em}
\noindent\textit{Note}: Results from cross-dataset experiments using the same model (Llama-2-7B) and unified hyperparameters at $p=0.10$. AGNEWS values are consistent with Table~\ref{tab:agnews_main}.
\end{table}

Table~\ref{tab:cross_dataset} reveals several key findings:

\begin{itemize}
    \item \textbf{Vanilla SFT vulnerability}: Without defense, harmful scores are extremely high (30.68--40.64\%), confirming the severity of the poisoning threat across all task types.
    \item \textbf{Classification tasks} (SST-2, AGNEWS): Both SafeGrad and SelGrad effectively reduce HS to $<$4\%. SelGrad achieves the lowest HS on AGNEWS (1.20\% vs SafeGrad 1.99\%), while SafeGrad achieves the lowest HS on SST-2 (1.99\% vs SelGrad 3.98\%). AsFT also provides effective defense on classification tasks (HS $<$4\%).
    \item \textbf{Reasoning task} (GSM8K): All defense methods show higher HS (5--7\%) due to task complexity. AsFT achieves the lowest HS (5.58\%), followed by SelGrad (6.37\%) and SafeGrad (7.17\%). However, AsFT's safety gain comes at a substantial FA cost (15.37\% vs SafeGrad 24.15\% and SelGrad 23.75\%), suggesting that AsFT's aggressive parameter anchoring disproportionately impairs reasoning capability.
\end{itemize}

%==================================================%
\section{Cross-Architecture Generalization}
\label{sec:cross_architecture}

This section evaluates SelGrad's generalization across different LLM architectures.

\subsection{Results Across Three Architectures}
\label{subsec:three_architectures}

Table~\ref{tab:cross_architecture} presents results across three model architectures on AGNEWS at $p=0.10$.

\begin{table}[htbp]
\centering
\caption{Cross-architecture generalization on AGNEWS ($p=0.10$)}
\label{tab:cross_architecture}
\renewcommand{\arraystretch}{1.2}
\begin{tabular}{l cc cc cc cc}
\toprule
& \multicolumn{2}{c}{\textbf{Vanilla SFT}} & \multicolumn{2}{c}{\textbf{AsFT}} & \multicolumn{2}{c}{\textbf{SafeGrad}} & \multicolumn{2}{c}{\textbf{SelGrad (Ours)}} \\
\cmidrule(lr){2-3} \cmidrule(lr){4-5} \cmidrule(lr){6-7} \cmidrule(lr){8-9}
\textbf{Model} & HS$\downarrow$ & FA$\uparrow$ & HS$\downarrow$ & FA$\uparrow$ & HS$\downarrow$ & FA$\uparrow$ & HS$\downarrow$ & FA$\uparrow$ \\
\midrule
Llama-2-7B & 30.68 & \cellcolor{bestcolor}86.70 & 2.39 & 82.20 & 1.99 & 86.20 & \cellcolor{bestcolor}1.20 & 85.60 \\
Llama-3.1-8B & 23.90 & 87.80 & 15.14 & \cellcolor{bestcolor}89.70 & 7.97 & 89.50 & \cellcolor{bestcolor}3.98 & 88.70 \\
Qwen-2-7B & 30.68 & 86.30 & \cellcolor{bestcolor}3.59 & 86.20 & 3.98 & \cellcolor{bestcolor}87.70 & \cellcolor{bestcolor}3.59 & 74.60 \\
\bottomrule
\end{tabular}

\vspace{0.3em}
\noindent\textit{Note}: All results on AGNEWS at $p=0.10$ with unified hyperparameters across three architectures. Values for Llama-2-7B are consistent with Table~\ref{tab:agnews_main}.
\end{table}

Table~\ref{tab:cross_architecture} reveals several architecture-dependent findings:

\paragraph{Analysis: Cross-Architecture Defense Effectiveness}
The results reveal important differences across architectures and defense methods:

\begin{enumerate}
    \item \textbf{SelGrad Outperforms SafeGrad on Llama-3.1-8B}: Contrary to earlier expectations, SelGrad achieves HS=3.98\% on Llama-3.1-8B, significantly outperforming SafeGrad (HS=7.97\%). This demonstrates that selective defense can be more effective than uniform gradient surgery even on architectures with improved pre-training.

    \item \textbf{AsFT Partial Failure on Llama-3.1-8B}: With corrected LR scheduling, AsFT achieves HS=15.14\% on Llama-3.1-8B---substantially better than Vanilla SFT (23.90\%) but still far worse than gradient-based defenses (SelGrad 3.98\%, SafeGrad 7.97\%). This indicates that AsFT's alignment-stage defense provides partial protection on Llama-3.1 but cannot match the effectiveness of gradient-level approaches. On Llama-2-7B and Qwen-2-7B, AsFT achieves competitive safety (HS=2.39\% and 3.59\%, respectively).

    \item \textbf{Qwen-2-7B SelGrad FA Degradation}: SelGrad achieves competitive safety on Qwen-2-7B (HS=3.59\%) but suffers a notable FA drop (74.60\% vs SafeGrad's 87.70\%). Qwen-2's unique tokenization creates architecture-specific loss patterns with higher overlap between benign and poisoned distributions, leading to more false positives and excessive suppression of task learning.
\end{enumerate}

\textbf{Recommendation}: For Qwen-family models, practitioners should apply architecture-specific threshold calibration to reduce false positives. For Llama-3.1, where AsFT provides only partial defense, gradient-based methods (SafeGrad or SelGrad) offer more reliable safety guarantees.

Figure~\ref{fig:cross_evaluation} visualizes the cross-dataset and cross-architecture generalization results, demonstrating SelGrad's safety-efficiency trade-off across diverse settings.

\begin{figure}[htbp]
    \centering
    \includegraphics[width=\textwidth]{figures/ch4_generalization.png}
    \caption{Cross-dataset and cross-architecture evaluation results at poisoning ratio $p=0.10$. (Left) Harmful Score comparison across three datasets: AsFT achieves the lowest HS on GSM8K (5.58\%), SelGrad achieves the lowest HS on AGNEWS (1.20\%), and SafeGrad achieves the lowest HS on SST-2 (1.99\%). Vanilla SFT shows high vulnerability across all datasets (30.68--40.64\%). (Right) Cross-architecture comparison across three architectures: SelGrad achieves the best HS on Llama-2-7B (1.20\%) and Llama-3.1-8B (3.98\% vs SafeGrad 7.97\%), while AsFT shows partial failure on Llama-3.1-8B (15.14\%).}
    \label{fig:cross_evaluation}
\end{figure}



%==================================================%
\section{Comparison with State-of-the-Art Methods}
\label{sec:sota_comparison}

This section provides comprehensive comparison with all baseline methods.

\subsection{Overall Performance Comparison}
\label{subsec:overall_comparison}

Table~\ref{tab:sota_comparison} presents a comprehensive comparison across all metrics.

\begin{table}[htbp]
\centering
\caption{Comparison of defense methods on AGNEWS ($p=0.10$, Llama-2-7B)}
\label{tab:sota_comparison}
\renewcommand{\arraystretch}{1.15}
\begin{tabular}{l cc cc}
\toprule
\textbf{Method} & \textbf{HS}$\downarrow$ & \textbf{FA}$\uparrow$ & \textbf{Time} & \textbf{VRAM} \\
\midrule
Vanilla SFT & 30.68 & \cellcolor{bestcolor}86.70 & \cellcolor{bestcolor}6.39m & \cellcolor{bestcolor}25.52G \\
AsFT & 2.39 & 82.20 & 24.36m & 33.40G \\
SafeGrad & 1.99 & 86.20 & 17.96m & 78.30G \\
\cellcolor{bestcolor}\textbf{SelGrad (Ours)} & \cellcolor{bestcolor}\textbf{1.20} & 85.60 & 8.97m & 35.86G \\
\bottomrule
\end{tabular}
\end{table}

Table~\ref{tab:sota_comparison} compares SelGrad against all baseline methods. SelGrad achieves the lowest harmful score (1.20\%) with competitive task accuracy, and 2$\times$ speedup over SafeGrad.

Figure~\ref{fig:comprehensive_comparison} provides a comprehensive visual comparison of all evaluated methods.

\begin{figure}[htbp]
    \centering
    \includegraphics[width=\textwidth]{figures/ch4_method_comparison.png}
    \caption{Comprehensive comparison with state-of-the-art methods. (Left) Heatmap showing Harmful Score across methods and poisoning ratios: darker colors indicate higher (worse) HS. SelGrad achieves the lowest or tied-lowest HS at moderate-to-high poisoning ratios while SafeGrad maintains slightly lower HS at low ratios. (Right) Efficiency comparison: SelGrad achieves 2$\times$ speedup over SafeGrad and 54\% VRAM reduction while maintaining competitive safety.}
    \label{fig:comprehensive_comparison}
\end{figure}

\subsection{Safety-Efficiency Trade-off Analysis}
\label{subsec:pareto_analysis}

SelGrad achieves a favorable position on the safety-efficiency trade-off frontier at moderate-to-high poisoning ratios ($p \geq 0.10$). Figure~\ref{fig:safety_accuracy_tradeoff} visualizes the safety-accuracy trade-off frontier.

\begin{figure}[htbp]
    \centering
    \includegraphics[width=\textwidth]{figures/ch4_safety_accuracy.png}
    \caption{Safety-accuracy trade-off across defense methods at $p=0.10$ on AGNEWS (Llama-2-7B). Each point represents a method's HS vs.\ FA. SelGrad achieves the lowest Harmful Score (1.20\%) with competitive Functional Accuracy (85.60\%). The optimal region (low HS, high FA) is highlighted in green shading.}
    \label{fig:safety_accuracy_tradeoff}
\end{figure}

SelGrad achieves the best safety-efficiency trade-off across multiple dimensions: it provides the strongest or tied-strongest safety (lowest or tied-lowest HS) at $p \geq 0.10$ with dramatically lower computational cost than SafeGrad, at a minor FA trade-off. This is enabled by selective defense application and LoRA-based memory optimization.

%==================================================%
\section{Ablation Study: Detection Signal Components}
\label{sec:ablation}

To validate that both identification criteria contribute meaningfully to SelGrad's effectiveness, we conduct an ablation study isolating each detection signal. Table~\ref{tab:ablation_detection} presents results on AGNEWS ($p=0.10$, Llama-2-7B) for four configurations: the full SelGrad system, each signal in isolation, and the SafeGrad baseline (uniform defense without selective identification).

\begin{table}[htbp]
\centering
\caption{Ablation study on detection signal components (AGNEWS, $p=0.10$, Llama-2-7B). Removing the early memorization signal causes catastrophic safety degradation, while both signals together achieve optimal performance.}
\label{tab:ablation_detection}
\begin{tabular}{lccc}
\toprule
\textbf{Variant} & \textbf{HS (\%)$\downarrow$} & \textbf{FA (\%)$\uparrow$} & \textbf{Time (min)} \\
\midrule
Full SelGrad (both signals) & \cellcolor{bestcolor}1.20 & 85.60 & 8.97 \\
Early memorization only ($\tau_{\text{low}}$) & 1.59 & 85.50 & 8.94 \\
SafeGrad (uniform defense) & 1.99 & \cellcolor{bestcolor}86.20 & 17.96 \\
High loss only ($\tau_{\text{high}}$) & 33.86 & 86.10 & \cellcolor{bestcolor}6.38 \\
\bottomrule
\end{tabular}
\end{table}

The results reveal a striking asymmetry between the two detection signals. The early memorization signal ($\tau_{\text{low}}$) is the dominant contributor to safety: using it alone achieves HS=$1.59\%$, close to the full system's $1.20\%$. In contrast, the high loss signal ($\tau_{\text{high}}$) alone yields HS=$33.86\%$---a catastrophic degradation comparable to undefended fine-tuning---because as the model memorizes poisoned samples, their loss decreases below $\tau_{\text{high}}$, and the adaptive threshold mechanism responds by lowering $\tau_{\text{high}}$ toward its floor (3.0), but this cannot keep pace with the rapid loss convergence of memorized samples. The early memorization signal prevents this failure mode by identifying poisoned samples precisely during the critical epochs 1--3 when backdoor associations are being formed. Combining both signals provides the best safety (HS=$1.20\%$), indicating that the high loss signal contributes a complementary $0.39$ percentage point improvement by catching poisoned samples that escape early detection in later epochs. From an efficiency perspective, the early memorization variant incurs nearly the same training time ($8.94$ min) as the full system ($8.97$ min), while the high loss variant is fastest ($6.38$ min) due to its lower detection rate triggering fewer gradient projection operations.

%==================================================%
\section{Discussion and Limitations}
\label{sec:limitations}

\subsection{Why Selective Defense Works}

Our results demonstrate that selective defense achieves the lowest harmful score (1.20\% vs SafeGrad 1.99\%) while maintaining competitive task accuracy. We identify three theoretical mechanisms that explain this effectiveness:

\paragraph{Gradient Interference Reduction}
SafeGrad applies uniform gradient surgery with alignment strength $\rho=1.0$ to \textit{all} training samples. At poison ratio $p=0.1$, this means 90\% benign samples receive unnecessary strong alignment, where alignment gradients $\nabla_\theta D_{KL}(f_{\theta_0} \| f_\theta)$ oppose task gradients (strongly conflicting).

SelGrad applies strong alignment ($\rho_p=0.25$) \textit{only} to identified suspicious samples (28.6\%), reducing gradient conflict by approximately 71\%. This allows benign samples to undergo free task learning, resulting in FA of 85.60\% compared to SafeGrad's 86.20\%---a minor trade-off ($-0.6\%$) for the significant safety improvement (HS from 1.99\% to 1.20\%).

\paragraph{Loss-based Patterns Capture Fundamental Attack Characteristics}
Our approach exploits two fundamental properties of backdoor poisoning:
\begin{itemize}
    \item \textbf{High Loss Pattern}: Backdoor triggers create unnatural input-output associations that the model struggles to fit initially, manifesting as elevated loss relative to benign samples.
    \item \textbf{Early Memorization Pattern}: Backdoors are easier to learn than task-specific patterns. In epochs 1--3, poisoned samples rapidly drop to loss $<2.0$ as simple trigger-target associations are memorized faster than genuine task semantics.
\end{itemize}

\paragraph{Sample Weighting Provides Layered Defense}
SelGrad combines pattern-driven selective surgery with sample weighting ($w_p=0.3$), creating a two-layer defense. Even if some poisoned samples are missed (false negatives), the combined weight reduction and surgery for identified samples significantly reduces backdoor learning.

\subsection{Applicability and Boundary Conditions}

While SelGrad demonstrates strong performance across diverse scenarios, we identify boundary conditions that inform deployment decisions:

\paragraph{Adversarially Optimized Attack Patterns}
SelGrad's loss-based identification assumes attackers use conventional poisoning strategies. Sophisticated adversaries with knowledge of our defense mechanism could potentially design triggers that mimic benign loss distributions more closely.

\textit{Theoretical Analysis}: Under adversarial optimization, the loss distribution separation $D_{\text{KL}}(P_b \| P_p)$ may decrease, reducing detection rates $\beta$ according to Theorem~\ref{thm:detection_bounds}. When $D_{\text{KL}} < 0.5$, we recommend falling back to uniform defense.

\textit{Defense-in-Depth Strategy}: We recommend combining SelGrad with complementary signals: (1) multi-modal identification integrating activation clustering and gradient norm analysis, (2) ensemble methods with multiple threshold configurations and voting mechanisms, (3) adaptive regularization with attention-based penalties.

\paragraph{Long-Context and Multi-Turn Tasks}
SelGrad's current thresholds are calibrated for short-context classification tasks. In long-context scenarios, benign samples naturally exhibit wider loss variance ($\sigma_b$ increases), affecting threshold calibration.

\textit{Extension Strategy}: Task-adaptive thresholds (automatically calibrate $\tau_{\text{high}}$ and $\tau_{\text{low}}$ using validation loss statistics), token-level detection for long contexts, and hybrid approaches combining with alignment-stage defenses for high-stakes applications.

\subsection{Architecture-Dependent Performance}

As demonstrated in Section~\ref{sec:cross_architecture}, SelGrad and competing methods exhibit notable variance across architectures:

\begin{itemize}
    \item \textbf{Strong Performance (Llama-2-7B)}: SelGrad achieves the lowest HS (1.20\%) among all methods on Llama-2-7B, with well-separated loss distributions enabling accurate detection.

    \item \textbf{Effective Defense (Llama-3.1-8B)}: SelGrad achieves HS=3.98\% on Llama-3.1-8B, outperforming SafeGrad (7.97\%) and AsFT (15.14\%). This demonstrates that selective defense can be more effective than both uniform gradient surgery and alignment-stage approaches on newer architectures.

    \item \textbf{Challenging (Qwen-2-7B)}: SelGrad achieves good safety (HS=3.59\%) but suffers significant FA degradation (74.60\% vs SafeGrad 87.70\%). Qwen-2's unique tokenization creates architecture-specific loss patterns with higher overlap between benign and poisoned distributions, leading to excessive false positives.

    \item \textbf{AsFT Architecture Sensitivity}: AsFT shows partial failure on Llama-3.1-8B (HS=15.14\%, better than Vanilla SFT at 23.90\% but far worse than gradient-based defenses), while achieving competitive safety on Llama-2-7B (HS=2.39\%) and Qwen-2-7B (HS=3.59\%). This indicates that AsFT's alignment-stage defense is sensitive to architectural differences, particularly for newer model families.
\end{itemize}

\textbf{Practical Deployment Guidelines}:
\begin{enumerate}
    \item \textbf{Llama-2 and Llama-3.1 families}: Deploy SelGrad with default parameters
    \item \textbf{Qwen family}: Apply architecture-specific threshold calibration to mitigate FA degradation, or use SafeGrad as fallback
    \item \textbf{New architectures}: Perform quick loss distribution analysis ($D_{\text{KL}}$ on 100-sample validation)---if $D_{\text{KL}} < 1.0$, consider hybrid defense strategies
\end{enumerate}

%==================================================%
\section{Summary}
\label{sec:results_summary}

This chapter presented comprehensive experimental evaluation of SelGrad, demonstrating:

\begin{enumerate}
    \item \textbf{Superior Safety}: SelGrad achieves the lowest or tied-lowest harmful score at moderate-to-high poisoning ratios ($p \geq 0.10$) on the primary benchmark.

    \item \textbf{Competitive Task Accuracy}: Free learning on benign samples maintains competitive functional accuracy while requiring dramatically less computation than uniform defense approaches.

    \item \textbf{Significant Efficiency Gains}: SelGrad achieves approximately 2$\times$ speedup and 54\% VRAM reduction compared to SafeGrad.

    \item \textbf{Robust Generalization}: Consistent performance across multiple datasets, model architectures, and poisoning ratios.

    \item \textbf{Favorable Trade-off}: Best balance of safety, accuracy, and efficiency at moderate-to-high poisoning ratios ($p \geq 0.10$).

    \item \textbf{Component Importance}: Ablation analysis (Table~\ref{tab:ablation_detection}) reveals that the early memorization signal is the dominant safety contributor (HS=$1.59\%$ alone vs.\ $33.86\%$ for high loss only), while combining both signals achieves optimal safety (HS=$1.20\%$).
\end{enumerate}

The next chapter discusses the implications of these results, limitations, and directions for future research.

\chapter{Conclusions}
\label{cha:conclusions}

This thesis addressed the critical challenge of preserving safety alignment in large language models during fine-tuning, particularly when training data may be contaminated with malicious exemplars. We identified a fundamental inefficiency pervading existing defense mechanisms: current approaches such as SafeGrad apply uniform gradient projection to all training samples, irrespective of whether they actually constitute a safety threat. This uniform treatment engenders unnecessary computational overhead and potential degradation of task-specific learning on benign samples.

To address this limitation, we proposed SelGrad (Selective Gradient Projection), a principled defense framework that achieves efficient safety alignment through intelligent sample-level discrimination. Our approach introduces two complementary identification mechanisms: (1) a loss-based criterion that detects harmful samples manifesting anomalously elevated initial loss values, and (2) an early memorization pattern detector that identifies samples exhibiting unusually precipitous loss convergence during early training iterations. Through combination of these indicators, SelGrad accurately discriminates potentially harmful samples from benign ones, enabling targeted application of gradient surgery exclusively to suspicious samples.

The core technical contribution of SelGrad resides in its selective gradient surgery mechanism. Rather than projecting all gradient updates onto the null space of the reference gradient---as done in uniform approaches---SelGrad partitions each training batch into benign and suspicious subsets. Benign samples receive standard gradient updates without modification, preserving their complete contribution to task learning. Only suspicious samples undergo gradient projection, wherein components conflicting with the safety-critical reference gradient are excised. This selective strategy fundamentally reduces computational cost while maintaining robust safety guarantees.

Additionally, we introduced a memory-efficient KL divergence computation methodology predicated on LoRA architecture. By computing reference model outputs through efficient weight reconstruction rather than maintaining a complete reference model copy in memory, our approach achieves approximately 54\% reduction in VRAM requirements relative to dual-model baseline implementations. This optimization renders SelGrad practically deployable on consumer-grade hardware with constrained GPU memory.

Experimental evaluation across multiple datasets and model architectures demonstrated the efficacy of our approach. Table~\ref{tab:conclusion_summary} summarizes the principal performance metrics comparing SelGrad against baseline methodologies on AGNEWS with Llama-2-7B at poisoning ratio $p=0.10$.

\begin{table}[htbp]
\centering
\caption{Summary of key performance metrics (AGNEWS, Llama-2-7B, $p=0.10$). SelGrad achieves the best safety-efficiency trade-off, with the lowest HS and dramatically reduced computational cost at a minor FA trade-off ($-0.6\%$).}
\label{tab:conclusion_summary}
\begin{tabular}{lcccc}
\toprule
\textbf{Method} & \textbf{HS (\%)$\downarrow$} & \textbf{FA (\%)$\uparrow$} & \textbf{Time (min)$\downarrow$} & \textbf{VRAM (GB)$\downarrow$} \\
\midrule
Vanilla SFT & 30.68 & 86.70 & 6.39 & 25.52 \\
AsFT & 2.39 & 82.20 & 24.36 & 33.40 \\
SafeGrad & 1.99 & 86.20 & 17.96 & 78.30 \\
\textbf{SelGrad (Ours)} & \textbf{1.20} & \textbf{85.60} & \textbf{8.97} & \textbf{35.86} \\
\midrule
\textbf{Improvement vs SafeGrad} & \textcolor{green!60!black}{-39.7\%} & \textcolor{red!60!black}{-0.6\%} & \textcolor{green!60!black}{-50.1\%} & \textcolor{green!60!black}{-54.2\%} \\
\bottomrule
\end{tabular}
\end{table}

SelGrad achieved a Harmful Score (HS) of 1.20\% at poisoning ratio $p=0.1$, compared to 30.68\% for undefended fine-tuning and 1.99\% for SafeGrad---representing a 39.7\% relative improvement. SelGrad completed training approximately 2$\times$ faster than SafeGrad (8.97 min versus 17.96 min) and required only 35.86 GB of VRAM compared to SafeGrad's 78.30 GB---a 54\% reduction. While SafeGrad achieved slightly higher FA (86.20\% versus SelGrad's 85.60\%), SelGrad's superior safety and dramatically lower computational cost make it the preferred choice for practical deployment.

Cross-dataset generalization experiments on SST-2 and GSM8K confirmed that SelGrad's detection mechanism transfers effectively across heterogeneous data distributions and task modalities. Cross-architecture evaluation on Llama-3.1-8B and Qwen-2-7B further validated the approach's applicability to diverse model families, with SelGrad outperforming SafeGrad on safety metrics in most configurations (e.g., HS of 3.98\% versus 7.97\% on Llama-3.1-8B). Ablation studies revealed that both identification criteria contribute meaningfully to detection accuracy, with their combination achieving optimal performance across varying poisoning ratios.

Notwithstanding these contributions, this study possesses several limitations warranting acknowledgment. First, our theoretical analysis presupposes that harmful and benign samples exhibit distinguishable loss characteristics---an assumption that may prove untenable for sophisticated adversarial attacks engineered to mimic benign sample behavior. Second, the detection thresholds ($\tau_{\text{high}}$ and $\tau_{\text{low}}$ absolute loss thresholds with adaptive adjustment, early memorization window) were determined empirically and may necessitate task-specific calibration for optimal performance. Third, while SelGrad demonstrates robust results on classification and instruction-following tasks, its efficacy on more complex generative tasks with extended output sequences remains to be comprehensively evaluated.

\section{Future Work}

Building on the findings and limitations of this thesis, we delineate concrete research directions with specific technical approaches:

\subsection*{Near-Term Extensions (6-12 months)}

\begin{enumerate}
    \item \textbf{Ensemble Detection with Gradient Signals:} Combine loss-based detection with gradient conflict frequency tracking. \textit{Technical approach}: Maintain per-sample conflict history $C_i = \sum_{t} \mathbf{1}[\mathbf{g}_{\text{task}}^{(i)} \cdot \mathbf{g}_{\text{align}} < 0]$ and flag samples where $C_i / T > 0.6$. This ensemble approach could improve robustness against loss-mimicking attacks where poisoned samples are designed to exhibit benign-like loss profiles.

    \item \textbf{Architecture-Adaptive Thresholds:} Address architecture-specific challenges through automatic threshold calibration. \textit{Technical approach}: Use validation-set loss statistics to set $\tau_{\text{high}} = \mu_{\text{val}} + k\sigma_{\text{val}}$ where $k$ is tuned via grid search on a small poison-labeled holdout (50 samples). This is particularly relevant for architectures like Qwen-2-7B where SelGrad's FA can be further improved through better-calibrated detection boundaries.

    \item \textbf{Memory-Efficient Scaling to 70B+ Models:} Current LoRA-based KL computation scales linearly with model size. \textit{Technical approach}: Implement activation checkpointing for reference forward pass and explore low-rank approximations of KL gradient ($\nabla D_{\text{KL}} \approx U\Sigma V^T$ with rank-16 truncation). Target: 70B model training on single RTX PRO 6000 96GB.
\end{enumerate}

\subsection*{Medium-Term Research (1-2 years)}

\begin{enumerate}
    \item \textbf{Multi-Modal SelGrad:} Extend to vision-language models (LLaVA, GPT-4V). \textit{Technical approach}: Compute separate loss signals for text and image modalities, apply cross-modal attention masking for samples flagged in either modality. Key challenge: visual backdoors may not manifest in text loss.

    \item \textbf{Federated Safe Fine-tuning:} Adapt for distributed training without centralizing data. \textit{Technical approach}: Each client computes local detection statistics, server aggregates via secure averaging to identify globally anomalous gradient contributions. Privacy guarantee: differential privacy ($\epsilon=1.0$) via gradient clipping and noise injection.

    \item \textbf{Certified Defense Bounds:} Derive PAC-Bayes bounds on HS under worst-case adaptive attacks. \textit{Technical approach}: Model attacker's optimization problem as finding poison samples that minimize detection probability while maximizing harm. Use robust optimization to derive certified upper bounds on HS given attacker's budget constraints.
\end{enumerate}

\subsection*{Long-Term Vision (2+ years)}

\begin{enumerate}
    \item \textbf{Universal Safety Primitives:} Develop model-agnostic safety modules that transfer across architectures. \textit{Research direction}: Learn a ``safety subspace'' from diverse aligned models and project all fine-tuning updates away from this subspace, regardless of detection.

    \item \textbf{Online Continual Safety:} Enable safe incremental updates without full retraining. \textit{Research direction}: Maintain a safety-critical parameter set $\mathcal{P}_{\text{safe}}$ identified via Fisher information, apply SelGrad-style protection only to $\mathcal{P}_{\text{safe}}$ during continual updates.
\end{enumerate}

In conclusion, SelGrad represents a significant advancement toward practical and efficient safety alignment for fine-tuned large language models. By recognizing that not all training samples necessitate defensive intervention, our selective approach achieves the dual objectives of computational efficiency and robust safety preservation. As LLM fine-tuning becomes increasingly democratized through open-source models and accessible computing resources, defense mechanisms such as SelGrad will assume a crucial role in ensuring that model customization does not compromise fundamental safety properties.

